% Options for packages loaded elsewhere
\PassOptionsToPackage{unicode}{hyperref}
\PassOptionsToPackage{hyphens}{url}
\documentclass[
  11pt,
  a4paper,
]{article}
\usepackage{xcolor}
\usepackage[margin=1in]{geometry}
\usepackage{amsmath,amssymb}
\setcounter{secnumdepth}{-\maxdimen} % remove section numbering
\usepackage{iftex}
\ifPDFTeX
  \usepackage[T1]{fontenc}
  \usepackage[utf8]{inputenc}
  \usepackage{textcomp} % provide euro and other symbols
\else % if luatex or xetex
  \usepackage{unicode-math} % this also loads fontspec
  \defaultfontfeatures{Scale=MatchLowercase}
  \defaultfontfeatures[\rmfamily]{Ligatures=TeX,Scale=1}
\fi
\usepackage{lmodern}
\ifPDFTeX\else
  % xetex/luatex font selection
  \setmainfont[ItalicFont=Songti TC,BoldItalicFont=Songti TC Bold]{Songti TC}
\fi
% Use upquote if available, for straight quotes in verbatim environments
\IfFileExists{upquote.sty}{\usepackage{upquote}}{}
\IfFileExists{microtype.sty}{% use microtype if available
  \usepackage[]{microtype}
  \UseMicrotypeSet[protrusion]{basicmath} % disable protrusion for tt fonts
}{}
\makeatletter
\@ifundefined{KOMAClassName}{% if non-KOMA class
  \IfFileExists{parskip.sty}{%
    \usepackage{parskip}
  }{% else
    \setlength{\parindent}{0pt}
    \setlength{\parskip}{6pt plus 2pt minus 1pt}}
}{% if KOMA class
  \KOMAoptions{parskip=half}}
\makeatother
\usepackage{longtable,booktabs,array}
\newcounter{none} % for unnumbered tables
\usepackage{calc} % for calculating minipage widths
% Correct order of tables after \paragraph or \subparagraph
\usepackage{etoolbox}
\makeatletter
\patchcmd\longtable{\par}{\if@noskipsec\mbox{}\fi\par}{}{}
\makeatother
% Allow footnotes in longtable head/foot
\IfFileExists{footnotehyper.sty}{\usepackage{footnotehyper}}{\usepackage{footnote}}
\makesavenoteenv{longtable}
% definitions for citeproc citations
\NewDocumentCommand\citeproctext{}{}
\NewDocumentCommand\citeproc{mm}{%
  \begingroup\def\citeproctext{#2}\cite{#1}\endgroup}
\makeatletter
 % allow citations to break across lines
 \let\@cite@ofmt\@firstofone
 % avoid brackets around text for \cite:
 \def\@biblabel#1{}
 \def\@cite#1#2{{#1\if@tempswa , #2\fi}}
\makeatother
\newlength{\cslhangindent}
\setlength{\cslhangindent}{1.5em}
\newlength{\csllabelwidth}
\setlength{\csllabelwidth}{3em}
\newenvironment{CSLReferences}[2] % #1 hanging-indent, #2 entry-spacing
 {\begin{list}{}{%
  \setlength{\itemindent}{0pt}
  \setlength{\leftmargin}{0pt}
  \setlength{\parsep}{0pt}
  % turn on hanging indent if param 1 is 1
  \ifodd #1
   \setlength{\leftmargin}{\cslhangindent}
   \setlength{\itemindent}{-1\cslhangindent}
  \fi
  % set entry spacing
  \setlength{\itemsep}{#2\baselineskip}}}
 {\end{list}}
\usepackage{calc}
\newcommand{\CSLBlock}[1]{\hfill\break\parbox[t]{\linewidth}{\strut\ignorespaces#1\strut}}
\newcommand{\CSLLeftMargin}[1]{\parbox[t]{\csllabelwidth}{\strut#1\strut}}
\newcommand{\CSLRightInline}[1]{\parbox[t]{\linewidth - \csllabelwidth}{\strut#1\strut}}
\newcommand{\CSLIndent}[1]{\hspace{\cslhangindent}#1}
\setlength{\emergencystretch}{3em} % prevent overfull lines
\providecommand{\tightlist}{%
  \setlength{\itemsep}{0pt}\setlength{\parskip}{0pt}}
% Journal-neutral XeLaTeX additions for the bilingual research draft.
\XeTeXlinebreaklocale "zh"
\XeTeXlinebreakskip = 0pt plus 1pt
\emergencystretch=2em
\setlength{\parindent}{1.5em}
\setlength{\parskip}{0.35em}
\AtBeginDocument{\hypersetup{
  linkcolor=blue!45!black,
  citecolor=blue!45!black,
  urlcolor=blue!55!black
}}
\usepackage{bookmark}
\IfFileExists{xurl.sty}{\usepackage{xurl}}{} % add URL line breaks if available
\urlstyle{same}
\hypersetup{
  pdftitle={Exact Smith Invariants and Affine Determinant Lines of Binary-Form Factorisation},
  pdfauthor={Anonymous},
  hidelinks,
  pdfcreator={LaTeX via pandoc}}

\title{Exact Smith Invariants and Affine Determinant Lines of
Binary-Form Factorisation}
\author{Anonymous}
\date{9 August 2026}

\begin{document}
\maketitle

\section{Abstract}\label{abstract}

For positive degrees \(\mathbf d\), put \(g=\gcd_i d_i\) and
\(\mathbf e=\mathbf d/g\), so that \(\mathbf e\) is primitive. Pairwise
resultants of labelled binary-form factors have characters
\(\chi_{ij}=[e_j\varepsilon_i+e_i\varepsilon_j]\) in
\(M=\mathbb Z^k/\mathbb Z\mathbf1\). Let \(W_E(\mathbf e)\) be the
matrix with these rows and let \(h(\mathbf e)\) be the positive gcd of
its maximal minors. We compute the full quotient by these characters.
For every prime \(p\), choose \(a\) with \(p\nmid e_a\). Then

\[
 \frac{M}{\langle\chi_{ij}\rangle}\otimes\mathbb Z_{(p)}
 \cong
 \frac{\mathbb Z_{(p)}}{
 (e_a-\sum_{i\ne a}e_i,\;2e_ie_j:i<j,\ i,j\ne a)}.
\]

This one-generator presentation gives every valuation of the final Smith
invariant, including the full \(2\)-adic case, a symmetric global
formula, and an \(O(k^2)\) gcd computation. It also shows that
spanning-tree minors already generate the full maximal-minor gcd.
Consequently the complete character matrix for the original degrees has
Smith form \(\operatorname{diag}(g,\ldots,g,g h(\mathbf e))\), with
\(k-2\) copies of \(g\).

We derive this lattice from an affine determinant-line theorem. Every
maximal minor of the differential of
\((A_1,\ldots,A_k)\mapsto\prod_iA_i\) equals, over \(\mathbb Z\), the
product of all pairwise resultants times the corresponding Plücker
coordinate of the product-one scaling torus, with an explicit
orientation. Multi-bordering contracts this tensor with a vertical torus
Jacobian; for semi-invariants that Jacobian is an integer character
determinant. On the universal pairwise-coprime open, pairwise resultants
generate all regular units modulo constants, so the complete lattice
defines the intrinsic unit-character residual group of that open, not a
classification of arbitrary semi-invariants. A square Laurent-unit
normalisation realises its index. More generally, for a square full-rank
semi-invariant character matrix \(W\) over an algebraically closed
field, the normalised factorisation map is finite locally free on a
dense target open of degree \(|\det W|D!/\prod_i d_i!\); its separable
and inseparable factors follow from the Smith entries of \(W\).

The square multiplication Jacobian, multifactor scalar resultant, signed
incidence support, arithmetic-matroid multiplicity, and labelled
root-partition count are classical inputs. The candidate increment is
the exact local/global calculation for the cross-weighted
diagonal-quotient character list, together with its factorisation
interpretation and the explicit integral affine Plücker lift. The
literature search is bounded, and the exact SymPy and FLINT checks are
producer-side regression evidence rather than independent reproduction.

\textbf{Keywords:} binary forms; resultants; Smith normal form;
arithmetic matroids; torus characters; determinant lines; factorisation
maps

\textbf{Parent work.} This paper is a separately versioned extension of
the unrefereed candidate \emph{Bordered Jacobian Foundations}; see the
\href{https://evidencepress.org/releases/bordered-jacobian-foundations/}{Evidence
Press parent page} and the immutable parent archive
\href{https://doi.org/10.5281/zenodo.21855302}{doi:10.5281/zenodo.21855302}.

\section{中文摘要}\label{ux4e2dux6587ux6458ux8981}

對正次數向量 \(\mathbf d\)，令 \(g=\gcd_i d_i\) 及
\(\mathbf e=\mathbf d/g\)，故 \(\mathbf e\)
為本原向量。標記二元齊次式各成對結式在特徵格
\(\mathbb Z^k/\mathbb Z\mathbf1\) 中給出交叉加權特徵
\(\chi_{ij}=[e_j\varepsilon_i+e_i\varepsilon_j]\)。令 \(W_E(\mathbf e)\)
為以這些特徵為列之矩陣，並以 \(h(\mathbf e)\)
表示其最大子式的正最大公因數。本文逐質數計算由全部 \(\chi_{ij}\)
生成之格的商：選取 \(p\nmid e_a\) 後，局部商為一生成元模， 其關係理想由
\(e_a-\sum_{i\ne a}e_i\) 與所有 \(2e_ie_j\;(i,j\ne a)\)
生成。此表示給出最後 Smith 不變因子的每個
\(p\)-進賦值（包括完整的二進情形）、對稱的整體公式，以及僅用生成樹
子式即可取得的最大子式公因數。因此未約化次數向量的 Smith 標準形為
\(\operatorname{diag}(g,\ldots,g,gh(\mathbf e))\)。

此格源自一個整數型仿射行列式線定理：多因子乘法微分的每一最大子式，
皆為全部成對結式之乘積乘上乘積為一縮放環面的相應 Plücker 座標，並有
明確定向。加入邊界函數後得到環面方向 Jacobian；對半不變量而言即為整數
特徵行列式。在通用的成對互素開集上，成對結式在常數之外生成所有正則
單位，故完整特徵格內在地描述該開集的「單位特徵」殘餘群，而非任意半不
變量的分類。對代數閉域上的滿秩方形特徵矩陣，正規化因式分解映射在稠密
目標開集上為有限局部自由，其次數為
\(|\det W|D!/\prod_i d_i!\)，可分與不可分部分由 Smith 不變因子決定。

方形乘法 Jacobian、多因子純量結式、有號關聯矩陣支撐、算術擬陣乘數與
標記根分割皆屬經典輸入。本文候選增量為上述交叉加權對角商特徵格的精確
局部與整體計算、其因式分解詮釋，以及明確的整數型仿射 Plücker 提升。
文獻檢索範圍有限；SymPy 與 FLINT 的精確檢查僅屬作者端回歸證據，並非
獨立重現。

\textbf{關鍵詞：} 二元齊次式；結式；行列式線；代數環面；Smith
標準形；因式分解映射

\section{1. Introduction}\label{introduction}

Let \(V_d\) be the affine coefficient space of binary forms of degree
\(d\). For two positive degrees \(r,s\), coefficient multiplication is
the morphism

\[
 m_{r,s}:V_r\times V_s\longrightarrow V_{r+s},\qquad (A,B)\longmapsto AB.
\]

Its source dimension exceeds its target dimension by one. The generic
kernel of its differential is the infinitesimal scaling direction
\((A,-B)\), induced by \((A,B)\mapsto(\lambda A,\lambda^{-1}B)\). The
parent result proves an exact integral refinement of this dimension
count: the complete signed maximal-minor vector of \(Dm_{r,s}\) is a
fixed sign times \(\operatorname{Res}(A,B)(A,-B)\)
(\citeproc{ref-AnonymousBorderedJacobian2026}{Anonymous 2026}). A border
covector pairs with this missing direction. When the border is the
differential of a bihomogeneous function, the pairing is its
scaling-degree difference. Appendix E includes a self-contained proof of
that base identity in the present coefficient and Sylvester conventions.

The rank-one vector identity is a determinant-line completion of the
classical square chart identities reviewed below. It is used here as the
base of an explicit all-factor construction, not as a standalone
priority claim.

The rank-one formula suggests a more general object. Given positive
degrees \(d_1,\ldots,d_k\), multiplication

\[
 m_{\mathbf d}:\prod_{i=1}^kV_{d_i}\longrightarrow V_D,
 \qquad D=\sum_i d_i,
\]

has relative dimension \(k-1\). The product-one torus

\[
 T=\{(\lambda_1,\ldots,\lambda_k)\in\mathbb G_m^k:
       \lambda_1\cdots\lambda_k=1\}
\]

acts within every fibre. Pairwise common roots are detected by the
resultants \(R_{ij}=\operatorname{Res}(A_i,A_j)\). The natural candidate
for the determinant line is therefore a product of the
pairwise-collision divisor and the Plücker line of the torus tangent
space.

Several exact shadows of this expectation are classical. For two
factors, coefficient-product equations have a square Jacobian equal, up
to the stated coordinate convention, to a Sylvester matrix and hence to
the resultant; this appears in Hensel-lifting, real-algebraic-geometry,
projective-tangent, and arithmetic settings
(\citeproc{ref-Artin2022}{Artin 2022, sec. 1.8}, Lemma 1.8.5;
\citeproc{ref-BasuPollackRoy2006}{Basu et al. 2006, sec. 4.2.1};
\citeproc{ref-Chipalkatti2003}{Chipalkatti 2003}, Lemma 5.7;
\citeproc{ref-BhargavaCremonaFisherGajovic2022}{Bhargava et al. 2022},
Lemma 2.6). For an arbitrary number of monic factors, Chaperon and López
de Medrano give a corank formula and prove that pairwise coprimality is
equivalent to local invertibility
(\citeproc{ref-ChaperonLopezDeMedrano2009}{Chaperon and López de Medrano
2009}, Theorem 3). Mahatab and Sampath prove that the determinant of the
canonical polynomial Chinese remainder map is the product of the
pairwise resultants (\citeproc{ref-MahatabSampath2015}{Mahatab and
Sampath 2015}, Theorem A.3). Generalized Sylvester and Koszul-complex
constructions also place the resultant in the relevant determinantal
ideals; Chardin proves divisibility of the maximal minors and a generic
gcd statement (\citeproc{ref-Chardin1993}{Chardin 1993, sec. IV}, Lemma
1 and Remark 1). Projective polynomial multiplication and its finite
root-partition fibres are classical
(\citeproc{ref-BreidingKohnSturmfels2024}{Breiding et al. 2024}), as are
factorisation quotients by scaling tori in the linear-factor case
(\citeproc{ref-Kurth1997}{Kurth 1997}). These antecedents determine the
claim boundary. We do not claim the square product-Jacobian identity,
the monic multifactor rank criterion, the scalar product
\(\prod R_{ij}\), projective root partitions, or incidence minors as
new.

The principal arithmetic object is the complete edge-character list. Put
\(g=\gcd_i d_i\) and \(\mathbf e=\mathbf d/g\). In the character lattice

\[
 M=X^*(T)=\mathbb Z^k/\mathbb Z\mathbf1,
\]

the primitive degrees \(\mathbf e\) give the rows of a matrix
\(W_E(\mathbf e)\):

\[
 \chi_{ij}=[e_j\varepsilon_i+e_i\varepsilon_j],\qquad
 C_{\mathbf e}=M/\langle\chi_{ij}:i<j\rangle.
\]

Let \(h(\mathbf e)\) be the positive gcd of the maximal minors of
\(W_E(\mathbf e)\). Section 7 proves that it equals \(|C_{\mathbf e}|\).

The main theorem computes \(C_{\mathbf e}\) locally at every prime by a
one-generator presentation. It follows that \(C_{\mathbf e}\) is cyclic,
that every \(v_p(|C_{\mathbf e}|)\) has a closed formula, and that the
final index has a symmetric factorisation-free expression. For
\(\mathbf d=g\mathbf e\), this yields the complete Smith form

\[
 \operatorname{diag}
 \bigl(\underbrace{g,\ldots,g}_{k-2},g h(\mathbf e)\bigr).
\]

The proof chain leading to this arithmetic calculation has four further
parts. First, an integral affine lift packages every maximal minor of
\(Dm_{\mathbf d}\) as \(\Delta_{\mathbf d}=\prod_{i<j}R_{ij}\) times the
matching torus Plücker coordinate, with a fixed orientation. Second,
Laplace expansion and Cauchy--Binet give the multi-border formula;
semi-invariant borders contribute \(\det(W)\prod_ag_a\). Third,
resultant characters reduce integrally to a cross-weighted all-negative
incidence problem, giving all square character minors and showing that
spanning trees alone determine the full-list index. Fourth, pairwise
resultants are proved to generate every regular unit modulo constants on
the universal pairwise-coprime open. Hence the complete lattice is
intrinsic in that regular-unit category, and its quotient is the
unit-character residual group of that open.

The contribution hierarchy is intentionally explicit.

{\def\LTcaptype{none} % do not increment counter
\begin{longtable}[]{@{}
  >{\raggedright\arraybackslash}p{(\linewidth - 2\tabcolsep) * \real{0.5000}}
  >{\raggedright\arraybackslash}p{(\linewidth - 2\tabcolsep) * \real{0.5000}}@{}}
\toprule\noalign{}
\begin{minipage}[b]{\linewidth}\raggedright
Layer
\end{minipage} & \begin{minipage}[b]{\linewidth}\raggedright
Status used in this paper
\end{minipage} \\
\midrule\noalign{}
\endhead
\bottomrule\noalign{}
\endlastfoot
Square multiplication Jacobian, multifactor scalar resultant, monic rank
criterion, and labelled root partitions & Classical antecedents \\
All-factor affine Plücker tensor and multi-border identity & Explicit
integral lift and formal contraction of classical scalar/kernel data;
candidate coordinate-level contribution \\
Signed support and powers of two in character minors & Classical
all-negative signed-incidence theory \\
Degree weights, appended balance row, and resultant interpretation &
Factorisation-specific specialisation \\
\(m(E)\), the maximal-minor gcd, and complex toric component count &
Standard arithmetic-matroid and toric-arrangement data \\
Local cyclic presentation, exact valuations, global formula, and scaled
Smith form for the cross-weighted diagonal-quotient list & Strongest
candidate theorem package in this paper \\
Generic degree \(|\det W|D!/\prod_i d_i!\) & Scheme-theoretic synthesis
of a labelled root-partition cover and a torus isogeny \\
\end{longtable}
}

For a square full-rank semi-invariant character matrix \(W\) over an
algebraically closed field, we prove after restriction to an explicit
dense target open that the normalised factorisation map is finite
locally free of degree

\[
 |\det W|\frac{D!}{\prod_i d_i!}.
\]

Smith form then separates its total, separable, and inseparable degrees.
A square Laurent basis of the full regular-unit character lattice always
realises the full-list index, but a square chart made from individual
resultants need not do so. The results do not classify arbitrary
polynomial semi-invariants, recognise level sets as affine spaces,
classify Keller maps, or imply a Hessian conjecture result. The exact
programs are producer-side regression checks; they are not independent
reproduction, formal verification, or peer review.

\section{2. Conventions and
antecedents}\label{conventions-and-antecedents}

\subsection{2.1 Coefficients and
resultants}\label{coefficients-and-resultants}

For a positive integer \(d\), write

\[
 A=\sum_{j=0}^{d}a_jX^{d-j}Y^j.
\]

After dehomogenising by \(X=1\), the coefficient vector
\((a_0,\ldots,a_d)\) is read in ascending powers of \(Y\). For forms of
degrees \(r\) and \(s\), the Sylvester matrix consists of the \(s\)
shifted rows of the first coefficient vector followed by the \(r\)
shifted rows of the second. In this convention,
\(\operatorname{Res}(X^r,Y^s)=1\). This choice fixes every sign below.

The classical interpretation of a resultant as a determinant goes back
to Cayley (\citeproc{ref-Cayley1848}{Cayley 1848}) and is developed
through modern determinant-of-complexes formalisms by Jouanolou,
Chardin, Demazure, and Gelfand--Kapranov--Zelevinsky
(\citeproc{ref-Jouanolou1991}{Jouanolou 1991};
\citeproc{ref-Chardin1993}{Chardin 1993};
\citeproc{ref-Demazure2012}{Demazure 2012};
\citeproc{ref-GelfandKapranovZelevinsky1994}{Gelfand et al. 1994}).
Jouanolou's work also contains a bordered-minor evaluation for generic
linear forms and extensive gradient and bordered-Macaulay identities
(\citeproc{ref-Jouanolou1991}{Jouanolou 1991, secs. 5.4.1--5.4.4};
\citeproc{ref-Jouanolou1995}{1995}; \citeproc{ref-Jouanolou1997}{1997}).
Those identities concern square elimination systems, resultant
complexes, or congruences modulo the resultant. The morphism studied
here differentiates multiplication with respect to factor coefficients
and is rectangular before borders are chosen.

Adjacent work on binary discriminants and resultants includes the
perfectness and equivariant resolutions of their Jacobian ideals
(\citeproc{ref-DAndreaChipalkatti2007}{D'Andrea and Chipalkatti 2007}).
That theory concerns the resultant hypersurface and its coefficient
derivatives. More directly, Chipalkatti proves that the differential of
projective two-factor binary-form multiplication is injective exactly
for coprime factors (\citeproc{ref-Chipalkatti2003}{Chipalkatti 2003},
Lemma 5.7), while Chaperon and López de Medrano analyse the corank and
singularities of multiplication of monic polynomials, including
arbitrary numbers of factors
(\citeproc{ref-ChaperonLopezDeMedrano2009}{Chaperon and López de Medrano
2009}, Theorems 1--4). These sources establish the rank-locus
background. They do not, in the expanded bounded corpus examined for
this project, state the universal non-monic all-factor
complementary-minor tensor with its integral orientation.

\subsection{2.2 Dimensions and the ordered torus
basis}\label{dimensions-and-the-ordered-torus-basis}

Fix \(k\ge2\) and positive degrees \(\mathbf d=(d_1,\ldots,d_k)\). Set

\[
 D=\sum_{i=1}^kd_i,\qquad n=D+1,
 \qquad t=k-1,qquad N=D+k=n+t.
\]

The differential \(M=Dm_{\mathbf d}\) is an \(n\times N\) matrix. Source
coordinates are ordered first by factor and then by coefficient. For
\(b=1,\ldots,t\), define the source tangent vector

\[
 \kappa_b=(0,\ldots,0,A_b,0,\ldots,0,-A_k),
\]

where the two nonzero blocks are the \(b\)-th and last factor blocks.
Let \(K\) be the \(t\times N\) matrix with rows \(\kappa_b\). Direct
differentiation gives

\[
 M K^{\mathsf T}=0.
\]

These rows correspond to the cocharacter basis \(e_b-e_k\) of the
product-one torus.

For a \(t\)-element subset \(I\subset\{0,\ldots,N-1\}\), let
\(M_{\widehat I}\) be the square matrix obtained by deleting the columns
in \(I\). Define

\[
 p_I(M)=(-1)^{\sum_{i\in I}i}\det M_{\widehat I}.
\]

The zero-based sign is part of the theorem, not shorthand for an
unspecified orientation. Write \(K_I\) for the \(t\times t\) submatrix
of \(K\) in the columns \(I\).

Finally, define the pairwise-collision product

\[
 \Delta_{\mathbf d}=\prod_{1\le i<j\le k}
 \operatorname{Res}(A_i,A_j).
\]

\subsection{2.3 Square and monic
shadows}\label{square-and-monic-shadows}

Fixing one leading coefficient and the product leading coefficient
removes a two-factor affine scaling direction and turns the remaining
product equations into a square Sylvester map. Artin identifies this
Jacobian directly (\citeproc{ref-Artin2022}{Artin 2022, sec. 1.8}, Lemma
1.8.5); Basu, Pollack, and Roy give the monic version
(\citeproc{ref-BasuPollackRoy2006}{Basu et al. 2006, sec. 4.2.1}), and
Bhargava, Cremona, Fisher, and Gajović prove over an arbitrary ring that
the monic-times-arbitrary chart has Jacobian equal to the resultant
(\citeproc{ref-BhargavaCremonaFisherGajovic2022}{Bhargava et al. 2022},
Lemma 2.6). The Stacks Project uses the same Sylvester Jacobian to
describe the coprime monic chart as étale
(\citeproc{ref-StacksProject2026}{The Stacks Project Authors 2026}, tag
00U0). For several monic factors, the product of pairwise resultants is
also supplied by the polynomial Chinese-remainder determinant
(\citeproc{ref-MahatabSampath2015}{Mahatab and Sampath 2015}, Theorem
A.3).

Thus neither the square Jacobian nor the multifactor scalar is a
contribution of this paper. Before any leading coefficient is fixed,
however, the differential has a \((k-1)\)-dimensional kernel. The
refinement recorded below retains the full complementary-minor tensor:
which maximal minor corresponds to which torus Plücker coordinate, with
what universal integral orientation. This lift is close to the classical
identities and is formal from a generic determinant-line viewpoint once
one chart minor and the kernel are known. The value of the treatment
below is its universal coordinate identity, integral orientation, and
compatibility with all complementary minors. No standalone priority
claim is made for the lift. The next section isolates the composition
rule used here.

\section{3. A determinant-line composition
lemma}\label{a-determinant-line-composition-lemma}

The induction factors \(k\)-fold multiplication through \((k-1)\)-fold
multiplication and one two-factor multiplication. The following
complementary-minor identity works over an arbitrary commutative ring,
so no division by a selected nonzero minor or generic-rank reduction is
needed.

\textbf{Lemma 3.1 (block-composition sign).} Let \(R\) be a commutative
ring and fix \(r,q\ge1\) and \(t_1\ge0\). Put

\[
 m=r+q-1,\qquad N=r+t_1+q,
\]

and suppose

\[
 P\in R^{r\times(r+t_1)},\qquad
 L\in R^{m\times(m+1)},\qquad
 S=\begin{pmatrix}P&0\\0&I_q\end{pmatrix}
 \in R^{(m+1)\times N}.
\]

Set \(M=LS\). Assume that \(K_P\in R^{t_1\times(r+t_1)}\),
\(\beta\in R^{m+1}\), and scalars \(c_P,c_L\in R\) satisfy

\[
 p_J(P)=c_P\det(K_P)_J\quad(|J|=t_1),
\]

\[
 L\beta=0,\qquad
 p_{\{\ell\}}(L)=c_L\beta_\ell\quad(0\le\ell\le m).
\]

Let \(\lambda\in R^N\) be a lift with \(S\lambda=\beta\), extend \(K_P\)
by zero columns to \(K_0=(K_P\;0)\), and order

\[
 K=\begin{pmatrix}K_0\\\lambda^{\mathsf T}\end{pmatrix}
\]

by the \(t_1\) inherited rows followed by the lifted outer row. Then for
every \((t_1+1)\)-element column set \(I\),

\[
 \boxed{
 p_I(M)=(-1)^{t_1}c_Pc_L\det K_I.
 }
\]

\textbf{Proof.} The block matrix \(S\) inherits the complementary-minor
identity

\[
 p_J(S)=c_P\det(K_0)_J\qquad(|J|=t_1).
\tag{3.1}
\]

Indeed, if \(J\) meets the identity block, both sides vanish by block
rank; if \(J\) lies in the \(P\)-source, the complementary square matrix
has diagonal blocks \(P_{\widehat J}\) and \(I_q\).

Fix \(I=\{i_0<\cdots<i_{t_1}\}\) and put \(C=I^c\), so \(|C|=m\).
Cauchy--Binet, with the intermediate rows in increasing order, gives

\[
 \det M_C
 =\sum_{\ell=0}^{m}\det L_{\widehat\ell}\det S_{\widehat\ell,C}
 =c_L\sum_{\ell=0}^{m}(-1)^\ell\beta_\ell
       \det S_{\widehat\ell,C}.
\tag{3.2}
\]

Expanding the \((m+1)\)-square matrix \([S_C\mid\beta]\) along its final
column yields

\[
 \sum_{\ell=0}^{m}(-1)^\ell\beta_\ell
       \det S_{\widehat\ell,C}
 =(-1)^m\det[S_C\mid\beta].
\tag{3.3}
\]

Since \(\beta=S\lambda\), multilinearity in the final column gives

\[
 \det[S_C\mid\beta]
 =\sum_{a=0}^{t_1}\lambda_{i_a}\det[S_C\mid S_{i_a}].
\tag{3.4}
\]

There are \(m-i_a+a\) elements of \(C\) larger than \(i_a\). Move the
last column \(S_{i_a}\) into increasing position and apply (3.1):

\[
 \begin{aligned}
 \det[S_C\mid S_{i_a}]
 &=(-1)^{m-i_a+a}\det S_{C\cup\{i_a\}}\\
 &=(-1)^{m+a+\sum I}c_P
   \det(K_0)_{I\setminus\{i_a\}}.
 \end{aligned}
\tag{3.5}
\]

The two occurrences of \(i_a\) cancel modulo two. On the other hand,
expansion of \(\det K_I\) along its last row is

\[
 \det K_I
 =\sum_{a=0}^{t_1}(-1)^{t_1+a}\lambda_{i_a}
   \det(K_0)_{I\setminus\{i_a\}}.
\tag{3.6}
\]

Equations (3.4)--(3.6) imply

\[
 \det[S_C\mid\beta]
 =(-1)^{m+\sum I+t_1}c_P\det K_I.
\tag{3.7}
\]

Substitution into (3.2)--(3.3) gives
\(\det M_C=(-1)^{\sum I+t_1}c_Pc_L\det K_I\). Multiplying by the
defining factor \((-1)^{\sum I}\) proves the result. \(\square\)

\textbf{Remark 3.2.} The factor \((-1)^{t_1}\) depends on the fixed
block order. It cannot be recovered from the unordered multiset of
degrees. This is visible already for the ordered triples \((1,1,2)\),
\((1,2,1)\), and \((2,1,1)\), whose final orientations are respectively
negative, positive, and negative.

\section{4. Affine Plücker lift of the multiplication
Jacobian}\label{affine-pluxfccker-lift-of-the-multiplication-jacobian}

Define

\[
E(\mathbf d)=
 \sum_{1\le i<j\le k}d_i(d_j+1)+\binom{k+1}{2}+1,
 \qquad
 \varepsilon_{\mathbf d}=(-1)^{E(\mathbf d)}.
\]

Since

\[
 \binom{k+1}{2}+1-\binom{k-1}{2}=2k,
\]

the same orientation is obtained by replacing the final two terms in
\(E(\mathbf d)\) with \(\binom{k-1}{2}\). This parity-equivalent form
makes the inductive contribution of the \(k-2\) inherited kernel
directions transparent.

\textbf{Theorem 4.1 (affine Plücker lift of the classical multiplication
Jacobian).} For every \(k\ge2\), every positive ordered degree vector
\(\mathbf d\), and every \((k-1)\)-element source-column set \(I\),

\[
 \boxed{
 p_I(Dm_{\mathbf d})
 =\varepsilon_{\mathbf d}\Delta_{\mathbf d}\det K_I.
 }
\]

The equality is an identity in the universal coefficient ring over
\(\mathbb Z\), and hence is compatible with arbitrary base change.

The square-chart scalar and the pairwise-coprime regularity criterion
are classical antecedents. The content recorded by Theorem 4.1 is their
simultaneous non-monic affine lift: it specifies every complementary
minor, the matching torus Plücker coordinate, and the orientation in one
universal integral identity. This exact packaging is used as the
structural bridge to the character-lattice calculation; it is not
presented as a new resultant determinant in isolation.

\textbf{Proof.} We induct on \(k\). For \(k=2\), the self-contained
two-factor argument in Appendix E gives

\[
 p_{\{i\}}(Dm_{d_1,d_2})
 =(-1)^{d_1(d_2+1)}
   \operatorname{Res}(A_1,A_2)K_i.
\]

The proposed exponent differs from \(d_1(d_2+1)\) by \(\binom32+1=4\),
so the signs agree.

Assume the theorem for \(k-1\) factors. Put

\[
 D'=d_1+\cdots+d_{k-1},\qquad B=A_1\cdots A_{k-1}.
\]

Factor multiplication as

\[
 \prod_{i=1}^kV_{d_i}
 \xrightarrow{\;m_{d_1,\ldots,d_{k-1}}\times\mathrm{id}\;}
 V_{D'}\times V_{d_k}
 \xrightarrow{\;m_{D',d_k}\;}V_D.
\]

The inner Plücker scalar is

\[
 \varepsilon_{d_1,\ldots,d_{k-1}}
 \prod_{1\le i<j<k}R_{ij}.
\]

The outer two-factor scalar is

\[
 (-1)^{D'(d_k+1)}\operatorname{Res}(B,A_k).
\]

Resultant multiplicativity in the fixed Sylvester convention gives the
exact identity

\[
 \operatorname{Res}(B,A_k)=\prod_{i<k}R_{ik}.
\]

No permutation of factor blocks is made, so no additional resultant sign
appears.

For the first \(k-1\) factors, use the inherited kernel rows

\[
 \eta_b=(0,\ldots,A_b,\ldots,-A_{k-1},0),
 \qquad b=1,\ldots,k-2.
\]

A lift of the outer kernel \((B,-A_k)\) is

\[
 \ell=(0,\ldots,0,A_{k-1},-A_k),
\]

because differentiating the inner product in the \(A_{k-1}\) scaling
direction sends \(A_{k-1}\) to \(B\). The standard torus rows satisfy

\[
 \kappa_b=\eta_b+\ell\quad(1\le b\le k-2),
 \qquad \kappa_{k-1}=\ell.
\]

The corresponding row-basis transformation is unitriangular and has
determinant one. Lemma 3.1, with \(t_1=k-2\), therefore gives the parity
recurrence

\[
 e_k\equiv e_{k-1}+D'(d_k+1)+(k-2)\pmod2.
\]

The closed exponent satisfies

\[
 E(d_1,\ldots,d_k)-E(d_1,\ldots,d_{k-1})
 =D'(d_k+1)+k,
\]

and \(k\equiv k-2\pmod2\). The scalar resultants combine to
\(\Delta_{\mathbf d}\), while the determinant-one row change gives the
displayed standard kernel minors. This completes the induction.
\(\square\)

For three factors the sign can be written more compactly as

\[
 \varepsilon_{d_1,d_2,d_3}
 =(-1)^{d_1(d_2+1)+(d_1+d_2)(d_3+1)+1}.
\]

Thus every two-column deletion satisfies

\[
 (-1)^{i+j}\det M_{\widehat{\{i,j\}}}
 =\varepsilon_{d_1,d_2,d_3}R_{12}R_{13}R_{23}\det K_{\{i,j\}}.
\]

This is the first higher-rank case of the induction used in this paper.

\subsection{4.1 Classical rank criterion and scheme-theoretic
refinement}\label{classical-rank-criterion-and-scheme-theoretic-refinement}

The coordinate identities yield the exact ideal equality

\[
 \boxed{
 I_{D+1}(Dm_{\mathbf d})
 =\Delta_{\mathbf d}I_{k-1}(K).
 }
\]

On the locus where every factor is nonzero, \(K\) has row rank \(k-1\).
On the pairwise-coprime locus, these rows are the entire kernel,
consistently with the known monic multifactor rank criterion
(\citeproc{ref-ChaperonLopezDeMedrano2009}{Chaperon and López de Medrano
2009}, Theorem 3). To see this in the present affine coordinates, let a
tangent vector \((H_1,\ldots,H_k)\) satisfy

\[
 \sum_{i=1}^kH_i\prod_{j\ne i}A_j=0.
\]

Reduce the equation modulo \(A_i\). Since \(A_i\) is coprime to
\(\prod_{j\ne i}A_j\), it divides \(H_i\). Degree equality gives
\(H_i=c_iA_i\) for a scalar \(c_i\). Substitution gives
\(\sum_i c_i=0\), precisely the tangent equation of the product-one
torus.

The classical criterion identifies pairwise root collision as the
generic rank-loss set. The ideal identity refines it in the universal
affine coefficient ring by retaining every maximal minor. Each generic
pairwise resultant is a primitive irreducible polynomial and the
different \(R_{ij}\) are nonassociate
(\citeproc{ref-Jouanolou1991}{Jouanolou 1991}). At the height-one prime
\((R_{ij})\), the \(K\)-minor obtained by choosing one fixed coefficient
column from each of the first \(k-1\) factor blocks is a product of
those coefficients and is not divisible by \(R_{ij}\); it is therefore a
unit in the height-one localisation. All other resultants are units
there, so the maximal-minor ideal has valuation one.

The exact statement after base change is slightly different. Let
\(A_{\mathbb Z}\) be the universal coefficient ring, let
\(A_{\mathbb Z}\to B\) be any ring map, and put \(J_B=I_{k-1}(K_B)\).
Polynomial base change gives unconditionally

\[
 I_{D+1}(Dm_{\mathbf d,B})=\Delta_{\mathbf d,B}J_B.
\tag{4.4}
\]

If \(B\) is a Noetherian normal domain and \(\mathfrak p\) is height
one, then \(B_{\mathfrak p}\) is a DVR and

\[
 v_{\mathfrak p}\!\left(I_{D+1}(Dm_{\mathbf d,B})\right)
 =\sum_{e\in E(K_k)}v_{\mathfrak p}(R_{e,B})
  +v_{\mathfrak p}(J_B),
\tag{4.5}
\]

where the valuation of an ideal is the minimum valuation of its
generators. Consequently the order is exactly one along \(\mathfrak p\)
provided that

\begin{enumerate}
\def\labelenumi{\arabic{enumi}.}
\tightlist
\item
  \(v_{\mathfrak p}(R_{ij,B})=1\);
\item
  every \(R_{e,B}\) with \(e\ne ij\) is a unit at \(\mathfrak p\); and
\item
  some maximal \(K_B\)-minor is a unit there.
\end{enumerate}

Equivalently, if a prime effective Cartier divisor has local equation
\(s\), the same conclusion holds when the pulled-back \(R_{ij}=us\) for
a unit \(u\), all other resultants and some maximal \(K\)-minor are
units at its generic point. A ramified pullback \(R_{ij}=us^m\)
contributes order \(m\), even when the underlying component is reduced.
Second collisions and torus-rank defects contribute the additional terms
visible in (4.5). Intersections and zero-factor strata can therefore
carry extra scheme structure through \(J_B\). The theorem does not
assert that the determinantal ideal is globally principal on the affine
source.

\textbf{Remark 4.2.} The first power of each pairwise resultant also
follows from multidegree bookkeeping once divisibility and the kernel
tensor are known. The induction is stronger: it fixes the complete
vector and its orientation without an irreducibility or generic-divisor
argument.

\section{5. Multi-border contraction}\label{multi-border-contraction}

Let \(g_1,\ldots,g_t\) be polynomials in the source coefficients, and
let \(G\) be their \(t\times N\) gradient matrix,

\[
 G_{aj}=\frac{\partial g_a}{\partial x_j}.
\]

The torus vector fields associated with the rows of \(K\) are

\[
 \delta_b=\sum_{j=0}^{N-1}K_{bj}\frac{\partial}{\partial x_j},
 \qquad b=1,\ldots,t.
\]

Thus \((KG^{\mathsf T})_{ba}=\delta_bg_a\).

\textbf{Corollary 5.1 (multi-border contraction).} With the
multiplication rows followed by the gradient rows in the order
\(g_1,\ldots,g_t\),

\[
 \boxed{
 \det\begin{pmatrix}M\\G\end{pmatrix}
 =(-1)^{nt+\binom t2}\varepsilon_{\mathbf d}
   \Delta_{\mathbf d}\det(\delta_bg_a)_{a,b=1}^t.
 }
\]

\textbf{Proof.} Expand the determinant along the first \(n\) rows. If
\(I\) is the set of \(t\) columns assigned to the gradient block, the
corresponding term is the product of \(\det M_{\widehat I}\) and
\(\det G_I\), times the Laplace sign. Substitution of
\(p_I(M)=(-1)^{\sum I}\det M_{\widehat I}\) cancels the part of the
Laplace sign depending on \(I\). The remaining parity is
\(nt+\binom t2\). Theorem 4.1 then gives

\[
 (-1)^{nt+\binom t2}\varepsilon_{\mathbf d}\Delta_{\mathbf d}
 \sum_{|I|=t}\det K_I\det G_I.
\]

Cauchy--Binet identifies the sum with \(\det(KG^{\mathsf T})\).
Transposing this final \(t\times t\) matrix does not change its
determinant, and its entries are the displayed vertical derivatives.
\(\square\)

Suppose now that each \(g_a\) is a \(T\)-semi-invariant. Relative to the
cocharacter basis \(e_b-e_k\), let its integer character be the \(a\)-th
row of \(W\). Then

\[
 \delta_bg_a=W_{ab}g_a,
\]

so the vertical matrix is \(\operatorname{diag}(g_1,\ldots,g_t)W\).

\textbf{Corollary 5.2 (character determinant).}

\[
 \boxed{
 \det D(m_{\mathbf d},g_1,\ldots,g_t)
 =(-1)^{nt+\binom t2}\varepsilon_{\mathbf d}
   \Delta_{\mathbf d}\det(W)\prod_{a=1}^tg_a.
 }
\]

For \(k=2\), \(W\) has one entry. If \(g\) has bidegree \((p,q)\), that
entry is \(p-q\). The degree-difference theorem is therefore exactly the
rank-one case of the character determinant.

Corollary 5.2 separates two obstructions. The collision factor
\(\Delta\) records additional kernel directions caused by common roots.
The integer \(\det W\) records whether the chosen borders recover the
generic torus directions. Even when \(\det W\ne0\), the collision
divisor remains.

\section{6. Resultant-character minors via weighted signless
incidence}\label{resultant-character-minors-via-weighted-signless-incidence}

\subsection{6.1 Edge characters}\label{edge-characters}

Under factor scaling, the pairwise resultant transforms as

\[
 R_{ij}(\lambda_iA_i,\lambda_jA_j)
 =\lambda_i^{d_j}\lambda_j^{d_i}R_{ij}(A_i,A_j).
\]

On the cocharacter lattice

\[
 X_*(T)=\{u=(u_1,\ldots,u_k)\in\mathbb Z^k:\sum_i u_i=0\},
\]

the edge \(ij\) therefore defines

\[
 \chi_{ij}(u)=d_ju_i+d_iu_j.
\]

Use coordinates \(u_1,\ldots,u_{k-1}\) and eliminate
\(u_k=-\sum_{b<k}u_b\). For a set \(H\) of \(k-1\) edges, let \(W_H\) be
the resulting square character matrix.

The graph associated with \(H\) need not be connected. Let \(\deg_H(i)\)
denote the valence of vertex \(i\). A connected component is \emph{odd
unicyclic} if its number of edges equals its number of vertices and its
unique cycle has odd length.

\subsection{6.2 Signed-graphic reduction and graph
minors}\label{signed-graphic-reduction-and-graph-minors}

Once the edge characters are written in these coordinates, their
maximal-minor calculation is a degree-weighted specialization of
classical signless-incidence minor theory. Up to transpose and
edge-orientation conventions, the relevant matrix is the signed
incidence matrix of the all-negative signed graph \(-H\). In that
language, balance is bipartiteness, the rank defect counts bipartite
components, and an unbalanced unicyclic component contributes a factor
two. These support, rank, and power-of-two facts are classical
(\citeproc{ref-Zaslavsky1982}{Zaslavsky 1982, sec. 7D}, Corollaries 7D.1
and 7D.3(d),(j), and sec.~8A, Lemmas 8A.2--8A.3;
\citeproc{ref-Zaslavsky1983Erratum}{Zaslavsky 1983}). We do not use the
uncorrected wording of Corollary 7D.3(g). The theorem below records the
factorisation-specific cross-degree transformation and the evaluation of
the appended degree row; it makes no claim to a new signed-incidence
classification.

\textbf{Theorem 6.1 (resultant-character minor formula).} The
determinant \(\det W_H\) can be nonzero only when \(H\) has exactly one
tree component and every other component is odd unicyclic. Suppose this
condition holds, let \(c\) be the number of components, and let
\(P\sqcup Q\) be the bipartition of the tree component. Then

\[
 \boxed{
 |\det W_H|
 =2^{c-1}
 \left|\sum_{i\in P}d_i-\sum_{j\in Q}d_j\right|
 \prod_{i=1}^kd_i^{\deg_H(i)-1}.
 }
\]

The right-hand side is integral. When the tree component is an isolated
vertex, use the bipartition \(P=\{i\},Q=\varnothing\); its degree
balance \(d_i\) cancels the apparent exponent \(-1\). Within the stated
component pattern, the determinant is nonzero exactly when the tree
component has nonzero bipartite degree balance.

\textbf{Proof.} Before eliminating \(u_k\), form the \((k-1)\times k\)
matrix \(C_H\) whose edge \(ij\) row has entry \(d_j\) in column \(i\),
entry \(d_i\) in column \(j\), and zero elsewhere. Append a last row
\((1,\ldots,1)\). Let \(S\in\operatorname{GL}_k(\mathbb Z)\) have
columns \(e_1-e_k,\ldots,e_{k-1}-e_k,e_k\). Its determinant is one, and

\[
 \begin{pmatrix}C_H\\ \mathbf 1^{\mathsf T}\end{pmatrix}S
 =\begin{pmatrix}W_H&*\\0&1\end{pmatrix}.
\]

Thus the absolute determinant of the augmented matrix is exactly
\(|\det W_H|\); in particular, eliminating the torus constraint
introduces no hidden factor of \(k\).

Multiply column \(i\) by \(d_i\). The edge \(ij\) row now has the common
factor \(d_id_j\), with a one in each endpoint column after that factor
is removed. The last row becomes \((d_1,\ldots,d_k)\). Comparing
determinants before cancelling any factor gives the integral identity

\[
 \left(\prod_i d_i\right)|\det W_H|
 =\left(\prod_i d_i^{\deg_H(i)}\right)
 \left|
 \det\begin{pmatrix}
 B_H\\ d_1&\cdots&d_k
 \end{pmatrix}
 \right|.
\tag{6.1}
\]

If every vertex is incident to an edge, cancellation gives

\[
 |\det W_H|
 =\left(\prod_{i=1}^kd_i^{\deg_H(i)-1}\right)
 \left|
 \det\begin{pmatrix}
 B_H\\ d_1&\cdots&d_k
 \end{pmatrix}
 \right|,
\]

where \(B_H\) is the unoriented edge--vertex incidence matrix. If the
unique tree component is an isolated vertex \(i\), the appended
determinant in (6.1) contains the factor \(d_i\); cancelling that factor
first gives the same displayed formula with no negative exponent. Thus
no division by a degree is used to establish integrality.

Reorder rows and columns by connected component. A tree component on
\(v\) vertices supplies \(v-1\) independent incidence rows. Appending
the restricted degree row produces determinant, up to sign,

\[
 \sum_{i\in P}d_i-\sum_{j\in Q}d_j,
\]

because the signed bipartition vector spans the kernel of its incidence
rows. An odd-unicyclic component has a square incidence matrix of
absolute determinant two. An even-unicyclic component is singular. A
component with more than one independent cycle produces row dependence
in a square maximal minor, while two tree components leave more than one
missing incidence row for the single degree border. Since \(H\) has
\(k-1\) edges, nonzero determinant therefore requires one tree component
and only odd-unicyclic remaining components. Multiplying the component
determinants proves the formula. \(\square\)

The incidence facts used in the last paragraph are classical. Grossman,
Kulkarni, and Schochetman determine the minors and Smith normal form of
unoriented graph incidence matrices
(\citeproc{ref-GrossmanKulkarniSchochetman1995}{Grossman et al. 1995}).
The odd-unicyclic invertibility criterion also appears in later inverse
formulas (\citeproc{ref-HessertMallik2023}{Hessert and Mallik 2023}).
The factorisation-specific content is the weighted transformation from
resultant characters to that incidence matrix.

\textbf{Corollary 6.2 (tree formula).} If \(H=T\) is a spanning tree,
then

\[
 \boxed{
 |\det W_T|
 =\left|\sum_{i\in P}d_i-\sum_{j\in Q}d_j\right|
  \prod_i d_i^{\deg_T(i)-1}.
 }
\]

For a star centred at vertex \(k\),

\[
 |\det W_T|
 =d_k^{k-2}\left|d_k-\sum_{i<k}d_i\right|.
\]

For three equal linear factors, every star has determinant one. Equal
degrees therefore do not force degeneration when the torus rank is
greater than one. This contrasts with the two-factor resultant
character, whose single weight is the degree difference.

\subsection{6.3 A three-linear-factor
example}\label{a-three-linear-factor-example}

Let \(d_1=d_2=d_3=1\) and write \(A_i=x_iX+y_iY\). If

\[
 A_jA_\ell=p_iX^2+q_iXY+r_iY^2,
 \qquad \{i,j,\ell\}=\{1,2,3\},
\]

then, in block-major source order \((x_1,y_1,x_2,y_2,x_3,y_3)\), the
multiplication differential and ordered torus kernel are

\[
 M=
 \begin{pmatrix}
 p_1&0&p_2&0&p_3&0\\
 q_1&p_1&q_2&p_2&q_3&p_3\\
 r_1&q_1&r_2&q_2&r_3&q_3\\
 0&r_1&0&r_2&0&r_3
 \end{pmatrix},
 \qquad
 K=
 \begin{pmatrix}
 x_1&y_1&0&0&-x_3&-y_3\\
 0&0&x_2&y_2&-x_3&-y_3
 \end{pmatrix}.
\]

Here \(\varepsilon_{(1,1,1)}=-1\). For example, deleting the \(x_1\)-
and \(x_2\)-columns gives

\[
 p_{\{0,2\}}(M)
 =-R_{12}R_{13}R_{23}\det
 \begin{pmatrix}x_1&0\\0&x_2\end{pmatrix}
 =-x_1x_2R_{12}R_{13}R_{23}.
\]

Choose the path normalising functions \((R_{12},R_{23})\). In the
cocharacter basis \((e_1-e_3,e_2-e_3)\), their character matrix is

\[
 W_{\{12,23\}}=
 \begin{pmatrix}1&1\\-1&0\end{pmatrix},
 \qquad \det W_{\{12,23\}}=1.
\]

The full edge-character matrix is

\[
 W_E=\begin{pmatrix}1&1\\0&-1\\-1&0\end{pmatrix},
\]

which has Smith form \(\operatorname{diag}(1,1)\). Thus two resultant
normalising functions remove the two-dimensional scaling torus without a
residual finite group scheme. The collision factor
\(R_{12}R_{13}R_{23}\) nevertheless remains in the bordered Jacobian.
Replacing the degrees by \((1,1,2)\) changes the same path matrix to
\(\left(\begin{smallmatrix}1&1\\-1&1\end{smallmatrix}\right)\), of
determinant two; the corresponding residual group scheme is \(\mu_2\),
nonreduced in characteristic two.

\section{7. Exact Smith invariants of the unit-character
lattice}\label{exact-smith-invariants-of-the-unit-character-lattice}

\subsection{7.1 Regular units and arithmetic-matroid
multiplicity}\label{regular-units-and-arithmetic-matroid-multiplicity}

Let \(A_{\mathbb Z}\) be the universal coefficient ring and put

\[
 U=D(\Delta_{\mathbf d})\subset\operatorname{Spec}A_{\mathbb Z}.
\]

The pairwise resultants are primitive irreducibles and are pairwise
nonassociate in the UFD \(A_{\mathbb Z}\). Hence

\[
 \boxed{
 \Gamma(U,\mathcal O_U)^\times
 =\left\{\varepsilon\prod_{i<j}R_{ij}^{n_{ij}}:
 \varepsilon\in\{\pm1\},\ n_{ij}\in\mathbb Z\right\}.
 }
\tag{7.1}
\]

Indeed, if \(f/\Delta^N\) and \(g/\Delta^M\) are inverse, then
\(fg=\Delta^{N+M}\), and unique factorisation forces every irreducible
factor of \(f\) and \(g\) to be one of the \(R_{ij}\). Valuations at the
distinct height-one primes \((R_{ij})\) give uniqueness. Over a field,
the same proof replaces \(\{\pm1\}\) by the field's multiplicative
group.

Thus all regular units on \(U\) are semi-invariant Laurent resultant
monomials. Their characters form exactly the row lattice \(L\) of the
complete edge-character matrix \(W_E\). This is the intrinsic character
lattice of regular units on the fixed pairwise-coprime open \(U\). It is
not the lattice of every polynomial semi-invariant, and it can enlarge
after deleting further divisors. We call

\[
 D\bigl(X^*(T)/L\bigr)
\]

the \textbf{unit-character residual group of \(U\)}.

Set

\[
 g=\gcd(d_1,\ldots,d_k),\qquad e_i=d_i/g,
 \qquad M=X^*(T)=\mathbb Z^k/\mathbb Z\mathbf1.
\]

For the primitive vector \(\mathbf e\), let

\[
 \chi_{ij}=[e_j\varepsilon_i+e_i\varepsilon_j]\in M,
 \qquad C_{\mathbf e}=M/\langle\chi_{ij}:i<j\rangle,
\]

and define

\[
 h(\mathbf e)=
 \gcd_{\substack{H\subset E(K_k)\\|H|=k-1}}
 |\det W_H(\mathbf e)|.
\tag{7.2}
\]

This gcd is standard representable arithmetic-matroid data. For a
sublist \(S\) of the primitive edge characters, put

\[
 m(S)=
 [M\cap\operatorname{span}_{\mathbb R}(S):\langle S\rangle_{\mathbb Z}].
\]

Then \(h(\mathbf e)=m(E)\), while a full-rank square sublist has
\(m(H)=|\det W_H|\). The equality between \(m(E)\) and the gcd of basis
multiplicities is the standard GCD rule (\citeproc{ref-Moci2012}{Moci
2012, sec. 2.2}; \citeproc{ref-DAdderioMoci2013}{D'Adderio and Moci
2013, secs. 1.4--1.5}). For the original list \(\mathbf d=g\mathbf e\),
the full-list multiplicity is \(g^{k-1}h(\mathbf e)\), not
\(h(\mathbf e)\).

\subsection{7.2 A one-generator local
presentation}\label{a-one-generator-local-presentation}

Write \(x_i\in M\) for the image of \(\varepsilon_i\), so that
\(\sum_i x_i=0\). The edge relations are

\[
 q_{ij}=e_jx_i+e_ix_j.
\]

For a prime \(p\), let \(R_p=\mathbb Z_{(p)}\) and use the convention
\(v_p(0)=\infty\).

\textbf{Theorem 7.1 (exact local Smith presentation).} Let \(k\ge3\),
let \(\mathbf e\) be primitive, and choose an index \(a\) with
\(p\nmid e_a\). Put

\[
 D_a=e_a-\sum_{i\ne a}e_i
\]

and

\[
 I_{p,a}=
 (D_a,\;2e_ie_j:i<j,\ i,j\ne a)R_p.
\]

Then \(1\mapsto x_a\) induces an isomorphism

\[
 \boxed{C_{\mathbf e}\otimes R_p\cong R_p/I_{p,a}.}
\tag{7.3}
\]

In particular, the presentation is valid for \(p=2\), for \(k=3\), and
when \(D_a=0\).

\textbf{Proof.} Primitivity supplies a \(p\)-unit \(e_a\). Over \(R_p\),
each star relation

\[
 e_ix_a+e_ax_i=0
\]

eliminates \(x_i=-(e_i/e_a)x_a\). The ambient relation becomes
\((D_a/e_a)x_a=0\), which is equivalent to \(D_ax_a=0\). Each remaining
edge \(ij\), with \(i,j\ne a\), becomes

\[
 -\frac{2e_ie_j}{e_a}x_a=0,
\]

equivalent to \(2e_ie_jx_a=0\). These are all generators and relations,
so the Tietze elimination proves the module isomorphism, not merely an
equality of orders. No division by \(2\), \(D_a\), or a nonunit degree
occurs. \(\square\)

The ideal in (7.3) is independent of the admissible pivot: two such
presentations are the annihilator of the same canonical local module,
and the edge between the pivots changes the displayed cyclic generator
by a unit.

\textbf{Corollary 7.2 (all local valuations).} For any \(p\)-unit pivot
\(a\),

\[
 \boxed{
 v_p(h(\mathbf e))=
 m_p(\mathbf e):=
 \min\left(
 v_p(D_a),
 \min_{\substack{i<j\\i,j\ne a}}
 [v_p(2)+v_p(e_i)+v_p(e_j)]
 \right).
 }
\tag{7.4}
\]

Moreover,

\[
 C_{\mathbf e}\otimes R_p\cong R_p/(p^{m_p(\mathbf e)}),
 \qquad
 C_{\mathbf e}\cong\mathbb Z/h(\mathbf e)\mathbb Z.
\tag{7.5}
\]

\textbf{Proof.} Every ideal in the DVR \(R_p\) is generated by an
element of least valuation, which gives the first isomorphism. The
zeroth Fitting ideal of the local cokernel is both \(I_{p,a}\) and the
localisation of the ideal generated by all maximal minors. Its positive
generator over \(\mathbb Z\) is \(h(\mathbf e)\), proving (7.4). The
local presentation is torsion, hence
\(C_{\mathbf e}\otimes\mathbb Q=0\); the finitely generated abelian
group \(C_{\mathbf e}\) is therefore finite. Its primary localisations
are cyclic, and their direct sum has pairwise coprime orders, so the
global group is cyclic of order \(h\). \(\square\)

This quotient module is the full-set module of the represented matroid
over \(\mathbb Z\). Arithmetic multiplicity remembers only its order;
the matroid-over-a-ring and DVR frameworks retain its isomorphism type
and local invariant sequence (\citeproc{ref-FinkMoci2016}{Fink and Moci
2016}, Definition 2.1, eq. (2.1), secs. 5 and 6.1). Theorem 7.1 is an
explicit computation of that standard object for this cross-weighted
factorisation list.

\subsection{7.3 Closed formula and tree-gcd
equality}\label{closed-formula-and-tree-gcd-equality}

For an odd prime \(p\), (7.4) is nonzero precisely when exactly two
degrees, say \(e_a,e_b\), are \(p\)-units and \(e_a\equiv e_b\pmod p\).
In that case

\[
 m_p(\mathbf e)=
 \min\left(v_p(e_a-e_b),
            \min_{c\notin\{a,b\}}v_p(e_c)\right).
\tag{7.6}
\]

For \(p=2\), let \(s\) be the number of odd \(e_i\). Primitivity gives
\(s\ge1\), and

\[
 m_2(\mathbf e)=
 \begin{cases}
 0,&s\text{ odd},\\[1mm]
 1,&s\ge4\text{ and }s\text{ even},\\[1mm]
 \min\left(v_2(D_a),1+\min_{c\notin\{a,b\}}v_2(e_c)\right),&s=2,
 \end{cases}
\tag{7.7}
\]

where \(a,b\) are the two odd indices in the last line. Swapping \(a\)
and \(b\) does not change the truncated valuation: if
\(E=\sum_{c\notin\{a,b\}}e_c\) and \(q=\min_cv_2(e_c)\), then
\(D_a+D_b=-2E\) is divisible by \(2^{q+1}\).

The primewise description assembles without factoring a maximal minor.
Define

\[
 G_{ab}=\gcd\{e_c:c\notin\{a,b\}\},\qquad
 Q_{ab}=\gcd(|e_a-e_b|,G_{ab}),
\]

and let \(Q_{ab}^{\mathrm{odd}}=Q_{ab}/2^{v_2(Q_{ab})}\). Put

\[
 \eta_2(\mathbf e)=
 \begin{cases}
 0,&s\text{ odd},\\
 1,&s\ge4\text{ and }s\text{ even},\\
 \min\bigl(v_2(D_a),1+v_2(G_{ab})\bigr),&s=2.
 \end{cases}
\]

\textbf{Theorem 7.3 (closed global index).} One has

\[
 \boxed{
 h(\mathbf e)=2^{\eta_2(\mathbf e)}
 \prod_{1\le a<b\le k}Q_{ab}^{\mathrm{odd}}.
 }
\tag{7.8}
\]

The odd factors in the product are pairwise coprime. Consequently
\(h(\mathbf e)\) is computable with \(O(k^2)\) integer gcd operations.
For each fixed \(a\), prefix and suffix gcds of the other \(k-1\)
degrees compute every excluded-pair gcd \(G_{ab}\) in \(O(k)\) gcd
operations; repeating over \(a\) gives the stated bound. This is an
arithmetic-operation count, not a unit-cost bit-complexity claim.

\textbf{Proof.} If an odd prime divides \(Q_{ab}\), it divides every
degree outside \(\{a,b\}\) and divides \(e_a-e_b\). It cannot divide
\(e_a\) or \(e_b\), else it would divide every primitive degree. Thus
\(\{a,b\}\) is exactly the pair of unit indices, and (7.6) gives
\(v_p(h)=v_p(Q_{ab})\). Conversely every exceptional odd prime divides
that \(Q_{ab}\). The unit pair is unique, so no odd prime divides two
different \(Q_{ab}\). Formula (7.7) supplies the remaining exponent at
two. Equality follows prime by prime. \(\square\)

The full-list gcd is even more economical than its definition suggests.

\textbf{Theorem 7.4 (spanning trees determine the full index).} If

\[
 \tau(\mathbf e)=
 \gcd_{T\text{ spanning tree of }K_k}|\det W_T(\mathbf e)|,
\]

then \(\tau(\mathbf e)=h(\mathbf e)\). After localising at \(p\), choose
any \(p\)-unit pivot \(a\). The star centred at \(a\) and its associated
\(\binom{k-1}{2}\) bent stars generate the full maximal-minor ideal.

\textbf{Proof.} By the spanning-tree determinant formula in Corollary
6.2, after choosing a \(p\)-unit vertex \(a\), the star centred at \(a\)
has determinant, up to sign,

\[
 e_a^{k-2}D_a,
\]

so it generates \(D_a\) up to a unit in \(R_p\). For distinct
\(i,j\ne a\), remove \(aj\) from the star and add \(ij\). The resulting
bent star has determinant

\[
 \pm e_a^{k-3}e_i(D_a+2e_j).
\]

Subtracting \(e_iD_a\) yields \(2e_ie_j\). Hence these tree minors
contain the ideal \(I_{p,a}\). Conversely every tree determinant is a
maximal minor, and the ideal of all maximal minors is
\(\operatorname{Fitt}_0(C_{\mathbf e}\otimes R_p)=I_{p,a}\). The two
ideals therefore agree at every prime, so their positive integer gcds
agree. \(\square\)

Odd-unicyclic bases remain relevant as individual arithmetic-matroid
multiplicities and square-chart determinants. Theorem 7.4 says only that
they do not lower the single full-list saturation index.

\subsection{7.4 Nonprimitive Smith form and residual group
scheme}\label{nonprimitive-smith-form-and-residual-group-scheme}

Because \(W_E(\mathbf d)=gW_E(\mathbf e)\), the primitive calculation
gives the complete answer.

\textbf{Corollary 7.5 (complete Smith form).} For \(k\ge3\),

\[
 \boxed{
 \operatorname{SNF}(W_E(\mathbf d))
 =\operatorname{diag}
 \bigl(\underbrace{g,\ldots,g}_{k-2},g h(\mathbf e)\bigr).
 }
\tag{7.9}
\]

Thus

\[
 X^*(T)/L
 \cong(\mathbb Z/g\mathbb Z)^{k-2}
 \oplus\mathbb Z/(g h(\mathbf e))\mathbb Z.
\tag{7.10}
\]

Indeed, the unimodular matrices that diagonalise the primitive matrix
put its scalar multiple into (7.9), which already satisfies Smith
divisibility.

Over \(\operatorname{Spec}\mathbb Z\), the unit-character residual group
is noncanonically

\[
 \boxed{
 D(X^*(T)/L)\cong
 \mu_g^{\,k-2}\times\mu_{g h(\mathbf e)}.
 }
\tag{7.11}
\]

It is finite locally free of rank \(g^{k-1}h(\mathbf e)\). Over a field
of characteristic \(p>0\), write \(g=p^\gamma g'\), \(h=p^m h'\), with
\(p\nmid g'h'\). Its connected diagonalizable factor and prime-to-\(p\)
étale direct factor are

\[
 \mu_{p^\gamma}^{\,k-2}\times\mu_{p^{\gamma+m}},
 \qquad
 \mu_{g'}^{\,k-2}\times\mu_{g'h'},
\tag{7.12}
\]

respectively. Their ranks are \(p^{\gamma(k-1)+m}\) and
\((g')^{k-1}h'\). Thus the scheme is étale exactly when \(p\nmid gh\),
and the number of geometric points in an algebraic closure is
\((g')^{k-1}h'\), not its total scheme rank in bad characteristic.

For comparison with toric arrangements, over \(\mathbb C\) the standard
identity \(m(S)=\#\pi_0(\bigcap_{\chi\in S}\ker\chi)\) gives exactly the
same full-list order (\citeproc{ref-Moci2012}{Moci 2012}, Lemma 5.4;
\citeproc{ref-DAdderioMoci2013}{D'Adderio and Moci 2013}, Lemma 4.1).
The group-scheme statement (7.11) is the characteristic-sensitive
refinement.

The rectangular all-edge character morphism

\[
 T\longrightarrow(\mathbb G_m)^{E(K_k)}
\]

has scheme-theoretic image \(D(L)\). The induced map \(T\to D(L)\) is
finite locally free of degree \([X^*(T):L]=g^{k-1}h\), with kernel
(7.11). On coordinate rings this is the free inclusion
\(\mathbb Z[L]\subset\mathbb Z[X^*(T)]\), with a basis given by coset
representatives. The degree is the degree onto this image, not a claim
about an unlabelled rectangular target containing every root-partition
branch.

\subsection{7.5 Diagnostic examples}\label{diagnostic-examples}

{\def\LTcaptype{none} % do not increment counter
\begin{longtable}[]{@{}
  >{\raggedright\arraybackslash}p{(\linewidth - 6\tabcolsep) * \real{0.2308}}
  >{\raggedleft\arraybackslash}p{(\linewidth - 6\tabcolsep) * \real{0.3077}}
  >{\raggedright\arraybackslash}p{(\linewidth - 6\tabcolsep) * \real{0.2308}}
  >{\raggedright\arraybackslash}p{(\linewidth - 6\tabcolsep) * \real{0.2308}}@{}}
\toprule\noalign{}
\begin{minipage}[b]{\linewidth}\raggedright
Degrees
\end{minipage} & \begin{minipage}[b]{\linewidth}\raggedleft
\(h(\mathbf e)\)
\end{minipage} & \begin{minipage}[b]{\linewidth}\raggedright
Smith form
\end{minipage} & \begin{minipage}[b]{\linewidth}\raggedright
Diagnostic point
\end{minipage} \\
\midrule\noalign{}
\endhead
\bottomrule\noalign{}
\endlastfoot
\((1,1,1)\) & \(1\) & \(\operatorname{diag}(1,1)\) & Saturated full
list \\
\((1,1,3)\) & \(3\) & \(\operatorname{diag}(1,3)\) & Odd bad prime
\(3\) \\
\((1,1,9)\) & \(9\) & \(\operatorname{diag}(1,9)\) & \(v_3(h)=2\) \\
\((1,1,4)\) & \(4\) & \(\operatorname{diag}(1,4)\) & \(v_2(h)=2\) \\
\((1,5,4)\) & \(8\) & \(\operatorname{diag}(1,8)\) & \(v_2(h)=3\) \\
\((1,2,2)\) & \(1\) & \(\operatorname{diag}(1,1)\) & Full list
saturated; square edge minors have absolute values \(3,2,2\) \\
\((1,1,1,1)\) & \(2\) & \(\operatorname{diag}(1,1,2)\) & Four odd
entries give exactly one factor of \(2\) \\
\((2,2,18)=2(1,1,9)\) & \(9\) & \(\operatorname{diag}(2,18)\) &
Nonprimitive content and odd valuation \\
\end{longtable}
}

The families

\[
 h(1,1,N)=N,
 \qquad
 h(1,1+2^q,2^q)=2^{q+1}\quad(q\ge1)
\tag{7.13}
\]

realise arbitrary higher odd and two-adic valuations. The vector
\((1,2,2)\) separates full-list saturation from square edge charts: no
two individual resultant characters form a unimodular basis, despite
\(h=1\).

\section{8. Generic-degree synthesis and
étaleness}\label{generic-degree-synthesis-and-uxe9taleness}

\subsection{8.1 Square semi-invariant
charts}\label{square-semi-invariant-charts}

Let \(\Bbbk\) be an algebraically closed field of characteristic
\(p\ge0\). A polynomial \(g\) is a \(T\)-semi-invariant of character
\(\chi\) if \(g(\lambda\cdot A)=\chi(\lambda)g(A)\). Suppose
\(g_1,\ldots,g_t\) are nonzero semi-invariants whose character matrix
\(W\in M_t(\mathbb Z)\) has nonzero determinant. Set

\[
 \Phi=(m_{\mathbf d},g_1,\ldots,g_t):
 \prod_iV_{d_i}\longrightarrow V_D\times\mathbb A^t
\]

and

\[
 X^\circ=D\!\left(\Delta_{\mathbf d}\prod_{a=1}^tg_a\right).
\]

Here the generic scheme degree means the degree of the induced finite
extension of function fields. The proof below constructs a dense target
open on which the map is finite locally free, so this degree is also the
constant scheme length of every fibre over that open.

The following statement combines the classical labelled root-partition
degree with the scheme order of the torus isogeny defined by \(W\). Its
role is to translate the character-lattice data into the geometry of a
normalised factorisation chart.

\textbf{Theorem 8.1 (finite normalised factorisation cover).} Fix the
labelled degrees, let \(\Bbbk\) be algebraically closed, and let
\(g_1,\ldots,g_t\) be nonzero semi-invariants with a square character
matrix \(W\) of nonzero determinant. There is a dense open
\(V_D^{\mathrm{good}}\subset V_D\) such that the restriction of \(\Phi\)
over \(V_D^{\mathrm{good}}\times(\mathbb G_m)^t\) is finite locally
free. In particular, \(\Phi\) is dominant and generically finite, with

\[
 \boxed{
 \deg_{\mathrm{gen}}\Phi
 =|\det W|\frac{D!}{\prod_i d_i!}.
 }
\]

If the Smith form of \(W\) is \(\operatorname{diag}(s_1,\ldots,s_t)\),
then the torus part of each generic fibre is a translate of

\[
 \ker\chi_W\cong\prod_{i=1}^t\mu_{s_i}.
\]

It has scheme order \(|\det W|\). If \(p>0\), write \(s_i=p^{a_i}s_i'\)
with \(p\nmid s_i'\). The generic separable degree is

\[
 \left(\prod_i s_i'\right)\frac{D!}{\prod_i d_i!},
\]

the number of geometric points in an algebraic closure of a generic
fibre is the same number, and the inseparable factor is
\(p^{\sum_i a_i}=p^{v_p(\det W)}\). On \(X^\circ\), the map is étale
exactly when \(p\nmid\det W\), with the characteristic-zero convention.

\textbf{Proof.} Let \(V_D^{\mathrm{sf}}\) be the open set of nonzero
binary forms with \(D\) distinct geometric roots, and define

\[
 Z^{\mathrm{sf}}
 =V_D^{\mathrm{sf}}\mathop{\times}_{\mathbb P(V_D)}
   \prod_{i=1}^k\mathbb P(V_{d_i}),
\]

where the second map is projective multiplication and the projective
source is restricted to the inverse image of the squarefree locus. The
projection \(\pi:Z^{\mathrm{sf}}\to V_D^{\mathrm{sf}}\) is the finite
étale root-partition cover: a point above \(F\in V_D^{\mathrm{sf}}\) is
a projective factorisation obtained by partitioning the roots of \(F\)
into labelled blocks of sizes \((d_1,\ldots,d_k)\). Its degree is

\[
 \nu=\frac{D!}{\prod_i d_i!}.
\tag{8.1}
\]

This is also the standard projective multiplication-fibre description
(\citeproc{ref-BreidingKohnSturmfels2024}{Breiding et al. 2024, sec.
10.2}, Proposition 10.7 and Corollary 10.8, p.~130). Finite étaleness
follows from the free symmetric-group action on ordered distinct-root
configurations; equivalently, projective multiplication is finite and
its differential is an isomorphism here because the affine differential
has only the \(T\)-tangent kernel
(\citeproc{ref-ChaperonLopezDeMedrano2009}{Chaperon and López de Medrano
2009}, Theorem 3). Scheme-theoretically, the free actions of the
relevant constant symmetric groups have scheme quotients and give fppf
torsors by (\citeproc{ref-StacksProject2026}{The Stacks Project Authors
2026}, tag 07S7); after base change along the torsor they are disjoint
unions of copies of the base. Finite locally free degree and étaleness
descend fpqc-locally (\citeproc{ref-StacksProject2026}{The Stacks
Project Authors 2026}, tags 02VO and 02VN).

Put \(U^{\mathrm{sf}}=m_{\mathbf d}^{-1}(V_D^{\mathrm{sf}})\). The map

\[
 q:U^{\mathrm{sf}}\longrightarrow Z^{\mathrm{sf}},\qquad
 (A_i)_i\longmapsto\left(\prod_iA_i,([A_i])_i\right),
\]

is a \(T\)-torsor: two affine tuples with the same product and
projective factors differ by a unique product-one scaling. For every
\(a\), the principal ideal \((g_a)\) is \(T\)-stable, so its zero
subscheme descends along the fpqc torsor \(q\) to a closed subscheme
\(Z_a\subset Z^{\mathrm{sf}}\) (\citeproc{ref-StacksProject2026}{The
Stacks Project Authors 2026}, tags 04TW and 023T). The dense open
\(U^{\mathrm{sf}}\) of the affine source cannot be contained in the zero
locus of the nonzero polynomial \(g_a\); faithful flatness therefore
makes \(Z_a\) proper.

Moreover, \(Z^{\mathrm{sf}}\) is irreducible: it is the nonzero
tautological-line pullback over the irreducible squarefree open in
\(\prod_i\mathbb P(V_{d_i})\). Hence \(\dim Z_a<\dim Z^{\mathrm{sf}}\).
Since \(\pi\) is finite, \(\pi(Z_a)\) is a proper closed subset of
\(V_D^{\mathrm{sf}}\). Thus

\[
 V_D^{\mathrm{good}}
 =V_D^{\mathrm{sf}}\setminus\bigcup_{a=1}^t\pi(Z_a)
\]

is dense and open, and every \(g_a\) is nonzero on every factorisation
branch above it.

Write

\[
 B_0=V_D^{\mathrm{good}},\qquad
 Z_0=\pi^{-1}(B_0),\qquad
 P_0=q^{-1}(Z_0),\qquad
 G=(\mathbb G_m)^t.
\]

Every \(g_a\) is invertible on \(P_0\), and equivariance gives a
morphism \(g:P_0\to G\). Define

\[
 \theta=(q,g):P_0\longrightarrow Z_0\times G.
\]

Then the restricted normalised factorisation map factors as

\[
 \Phi_0=(m_{\mathbf d},g):P_0
 \xrightarrow{\ \theta\ }Z_0\times G
 \xrightarrow{\ \pi\times\mathrm{id}_G\ }B_0\times G.
\tag{8.2}
\]

Base-change \(\theta\) by the fpqc cover
\(q\times\mathrm{id}_G:P_0\times G\to Z_0\times G\). The torsor identity
\(P_0\times_{Z_0}P_0\cong P_0\times T\) identifies the pullback with the
Cartesian square

\[
\begin{array}{ccc}
 P_0\times T&\xrightarrow{\ \beta\ }&P_0\times G\\[1mm]
 \downarrow{a}&&\downarrow{q\times\mathrm{id}_G}\\[1mm]
 P_0&\xrightarrow{\ \theta\ }&Z_0\times G,
\end{array}
\tag{8.3}
\]

where

\[
 a(x,\lambda)=\lambda x,
 \qquad
 \beta(x,\lambda)=\bigl(x,g(x)\chi_W(\lambda)\bigr).
\]

The square is Cartesian because two points in one \(q\)-fibre differ by
a unique \(\lambda\in T\), and \(g(\lambda x)=\chi_W(\lambda)g(x)\).
Smith normal form and torus automorphisms reduce \(\chi_W\) to

\[
 (z_1,\ldots,z_t)\longmapsto(z_1^{s_1},\ldots,z_t^{s_t}).
\tag{8.4}
\]

On coordinate rings each power map is free of rank \(s_i\), with basis
\(1,z_i,\ldots,z_i^{s_i-1}\). Hence \(\beta\) is finite locally free of
rank \(|\det W|\). Finite local freeness and constant rank descend
fpqc-locally (\citeproc{ref-StacksProject2026}{The Stacks Project
Authors 2026}, tag 02VO), so \(\theta\) has the same rank. The second
map in (8.2) is finite étale of rank \(\nu\); finite locally free
morphisms are stable under composition
(\citeproc{ref-StacksProject2026}{The Stacks Project Authors 2026}, tag
02K9). Thus \(\Phi_0\) is finite locally free of rank \(\nu|\det W|\),
proving total scheme degree without counting geometric points.

For (8.4), the function-field extension in coordinate \(i\) has total
degree \(s_i\), separable degree \(s_i'\), and inseparable degree
\(p^{a_i}\). The root-partition cover is étale. Composition therefore
gives the displayed total and separable degrees and the inseparable
factor \(p^{\sum_i a_i}\). Scheme-theoretically,
\(\ker\chi_W\cong\prod_i\mu_{s_i}\) has length \(\prod_i s_i\), while
its geometric-point count is \(\prod_i s_i'\). Finally, Corollary 5.2
gives on \(X^\circ\)

\[
 \det D\Phi
 =\pm\Delta_{\mathbf d}\det(W)\prod_ag_a.
\]

All factors except \(\det W\) are units there, so the Jacobian criterion
gives the asserted étaleness condition; equivalently, étaleness is
fpqc-local (\citeproc{ref-StacksProject2026}{The Stacks Project Authors
2026}, tag 02VN). \(\square\)

\subsection{8.2 Full-unit, edge, and polynomial-monomial
charts}\label{full-unit-edge-and-polynomial-monomial-charts}

Choose a \(\mathbb Z\)-basis \(\ell_1,\ldots,\ell_t\) of the complete
regular-unit character lattice \(L\). Expressing each \(\ell_a\) as an
integral combination of edge characters produces Laurent units

\[
 u_a=\prod_{i<j}R_{ij}^{z_{a,ij}}\in\Gamma(U,\mathcal O_U)^\times.
\]

Their square character matrix has determinant
\([X^*(T):L]=g^{k-1}h(\mathbf e)\). Theorem 8.1 therefore gives the
Laurent-unit normalisation the generic degree

\[
 g^{k-1}h(\mathbf e)\frac{D!}{\prod_i d_i!}.
\tag{8.5}
\]

This realises the full-list index on the pairwise-coprime open, but
negative exponents need not extend across its boundary. By contrast, a
set \(H\) of \(t\) individual resultants has torus degree
\(|\det W_H|=m(H)\), while nonnegative products of resultants correspond
to nonnegative integral combinations of the rows. Arbitrary polynomial
semi-invariants form a still larger category and may require deletion of
additional divisors before becoming units.

The primitive vector \((1,2,2)\) makes the distinction concrete. In edge
order \(12,13,23\),

\[
 W_E=
 \begin{pmatrix}
 2&1\\1&-1\\-2&0
 \end{pmatrix},
\]

whose maximal minors have absolute values \(3,2,2\). Hence \(h=1\), but
no pair of individual resultants is unimodular. Nevertheless the
nonnegative monomials

\[
 u_1=R_{12}R_{13}R_{23},\qquad u_2=R_{12}R_{23}
\]

have character matrix

\[
 \begin{pmatrix}1&1&1\\1&0&1\end{pmatrix}W_E=I_2.
\tag{8.6}
\]

Here \(D=5\) and \(\nu=5!/(1!2!2!)=30\). Thus:

{\def\LTcaptype{none} % do not increment counter
\begin{longtable}[]{@{}lrr@{}}
\toprule\noalign{}
Normalising functions & Torus degree & Total generic degree \\
\midrule\noalign{}
\endhead
\bottomrule\noalign{}
\endlastfoot
\((R_{12},R_{13})\) & \(3\) & \(90\) \\
\((R_{12},R_{23})\) & \(2\) & \(60\) \\
\((R_{13},R_{23})\) & \(2\) & \(60\) \\
\((u_1,u_2)\) & \(1\) & \(30\) \\
\end{longtable}
}

The last chart is étale on \(D(\Delta)\) in every characteristic. It
proves that full-list saturation need not be realised by an edge basis
and that a polynomial monomial chart can outperform every edge basis. It
does not prove that every saturated resultant-character semigroup has a
nonnegative unimodular basis.

For \(k=2\) and degrees \((1,2)\), a rank-one normalising semi-invariant
of unit character determinant gives generic degree \(3\), the number of
ways to choose one of three roots for the labelled linear factor. For
three equal linear factors, a tree of resultant normalising functions
can have \(|\det W|=1\), while the root-partition factor is \(3!=6\).
Neither observation produces a polynomial Keller map: global
affine-slice recognition is an additional requirement.

\section{9. Exact computational
evidence}\label{exact-computational-evidence}

The supplied programs test the formulas over exact integer polynomial
rings. They do not use floating point or random sampling in their
default tiers. Their purpose is to detect convention, sign, exponent,
and implementation errors in the stated identities.

The first multifactor implementation uses SymPy 1.14 and Berkowitz
determinants. It constructs the coefficient differential directly from
shifted products of complementary factors, constructs each pairwise
Sylvester resultant independently, and compares every Plücker coordinate
with the predicted torus minor. Its seven default degree vectors are

\[
 (1,1),\ (1,2),\ (1,1,1),\ (1,1,2),\ (1,2,1),\ (2,1,1),\
 (1,1,1,1).
\]

These cases comprise 143 symbolic Plücker coordinates. Five deliberate
mutations omit one resultant, square the collision product, reverse a
kernel row, reverse the recursive sign, or perturb the multiplication
differential. All are detected.

The second implementation uses python-flint 0.9, multivariate integer
polynomials, and fraction-free Bareiss determinants. It does not import
the SymPy implementation. Its default tier checks 134 symbolic Plücker
coordinates for three- and four-factor cases. It also computes the full
square bordered Jacobian for the normalising functions
\((R_{12},R_{23})\) at degrees \((1,1,1)\) and \((1,1,2)\), whose
character determinants are respectively one and two. Four structural
mutations are detected.

A third program operates on exact integer character matrices. Through
five vertices and degrees at most three, it checks 31,761 spanning-tree
formulas, 52,741 general graph-minor formulas, and 351
Smith-form/bad-prime cases. Four negative controls substitute total
degree for bipartite balance, omit the odd-cycle factor two, omit the
common gcd in the Smith form, or extrapolate the two-factor equal-degree
collapse to three factors.

A fourth Stage-4 program targets the strengthened Smith theorem. It
constructs the complete character matrix independently, computes its
Smith form over \(\mathbb Z\), compares the last invariant with the
symmetric formula, and checks the local annihilator valuation for every
admissible pivot and every prime in its bounded domains. A separate tier
enumerates actual spanning trees by Prüfer sequences and takes the gcd
of their determinants. The default run contains 42,727 local/global
Smith comparisons, 1,485 tree-gcd comparisons, and eight
displayed-example checks, for 44,220 positive comparisons. Five negative
controls detect omission of the factor two, replacement of valuations by
prime support, conflation of full-list saturation with an edge basis,
and omission of nonprimitive content.

{\def\LTcaptype{none} % do not increment counter
\begin{longtable}[]{@{}
  >{\raggedright\arraybackslash}p{(\linewidth - 8\tabcolsep) * \real{0.1765}}
  >{\raggedleft\arraybackslash}p{(\linewidth - 8\tabcolsep) * \real{0.2353}}
  >{\raggedleft\arraybackslash}p{(\linewidth - 8\tabcolsep) * \real{0.2353}}
  >{\raggedright\arraybackslash}p{(\linewidth - 8\tabcolsep) * \real{0.1765}}
  >{\raggedright\arraybackslash}p{(\linewidth - 8\tabcolsep) * \real{0.1765}}@{}}
\toprule\noalign{}
\begin{minipage}[b]{\linewidth}\raggedright
Tier
\end{minipage} & \begin{minipage}[b]{\linewidth}\raggedleft
Positive comparisons
\end{minipage} & \begin{minipage}[b]{\linewidth}\raggedleft
Negative controls
\end{minipage} & \begin{minipage}[b]{\linewidth}\raggedright
Script SHA-256 prefix
\end{minipage} & \begin{minipage}[b]{\linewidth}\raggedright
Receipt SHA-256 prefix
\end{minipage} \\
\midrule\noalign{}
\endhead
\bottomrule\noalign{}
\endlastfoot
SymPy multifactor & 143 coordinates & 5 & \texttt{3806a7c9b4bf} &
\texttt{27d4ea5a3e89} \\
FLINT multifactor/borders & 134 coordinates + 2 borders & 4 &
\texttt{e7de053c910d} & \texttt{f352edfd686e} \\
Character lattice & 84,853 comparisons & 4 & \texttt{b91d2c7b0f11} &
\texttt{5efd7944b92c} \\
Local Smith theorem & 44,220 comparisons & 5 & \texttt{09ce5ad767d8} &
\texttt{4d437bf2a6c8} \\
\end{longtable}
}

The package manifest and Stage-1 provenance record contain the complete
64-character digests; the table uses unambiguous prefixes for
legibility.

Across the four tiers there are 129,352 positive comparisons and 18
detected negative controls. Each default tier was run under ordinary
Python and \texttt{python\ -O}. The normal and optimized receipts are
byte-identical. This guard matters because proof checks implemented only
with Python \texttt{assert} statements can disappear under optimization.

The two polynomial programs use distinct expression engines and
determinant algorithms, which reduces the probability of a
backend-specific error. They still implement the same producer
specification and were written within the same project. Their agreement
is cross-implementation evidence, not independent reproduction. Exact
computation also does not prove novelty or replace the integral
argument.

\section{10. Antecedents, limitations, and the remaining
programme}\label{antecedents-limitations-and-the-remaining-programme}

\subsection{10.1 Contribution boundary}\label{contribution-boundary}

The square multiplication-Jacobian identity is a direct antecedent, not
a novelty claim. Artin identifies the two-factor product-equation
Jacobian with the transposed resultant matrix
(\citeproc{ref-Artin2022}{Artin 2022, sec. 1.8}, Lemma 1.8.5); Basu,
Pollack, and Roy state the monic version
(\citeproc{ref-BasuPollackRoy2006}{Basu et al. 2006, sec. 4.2.1}); and
Bhargava, Cremona, Fisher, and Gajović give monic-times-arbitrary
multiplication over an arbitrary ring
(\citeproc{ref-BhargavaCremonaFisherGajovic2022}{Bhargava et al. 2022},
Lemma 2.6). Chipalkatti supplies the corresponding projective tangent
injectivity criterion (\citeproc{ref-Chipalkatti2003}{Chipalkatti 2003},
Lemma 5.7). Chaperon and López de Medrano analyse the singularities of
polynomial multiplication and prove the arbitrary-factor monic corank
and pairwise-coprimality criterion
(\citeproc{ref-ChaperonLopezDeMedrano2009}{Chaperon and López de Medrano
2009}, Theorems 1 and 3).

The product of pairwise resultants as a multifactor scalar is likewise a
direct antecedent: Mahatab and Sampath's polynomial CRT theorem gives it
for monic factors (\citeproc{ref-MahatabSampath2015}{Mahatab and Sampath
2015}, Theorem A.3). Chardin's Koszul-complex treatment places the
resultant in generalized Sylvester determinantal ideals
(\citeproc{ref-Chardin1993}{Chardin 1993, sec. IV}, Lemma 1 and Remark
1). The determinant-of-complexes principle is classical
(\citeproc{ref-Jouanolou1991}{Jouanolou 1991};
\citeproc{ref-Demazure2012}{Demazure 2012};
\citeproc{ref-GelfandKapranovZelevinsky1994}{Gelfand et al. 1994}).

The projective multiplication map, finite factorisation fibres, and
labelled root partitions are likewise established background
(\citeproc{ref-BreidingKohnSturmfels2024}{Breiding et al. 2024}).
Kurth's use of ordered linear factors modulo a scaling torus and the
symmetric group is a close quotient-space antecedent
(\citeproc{ref-Kurth1997}{Kurth 1997}). The signed support, rank, and
powers of two in Theorem 6.1 are classical all-negative signed-incidence
facts (\citeproc{ref-Zaslavsky1982}{Zaslavsky 1982},
\citeproc{ref-Zaslavsky1983Erratum}{1983};
\citeproc{ref-GrossmanKulkarniSchochetman1995}{Grossman et al. 1995}).
The equal-weight vectors \(\varepsilon_i+\varepsilon_j\) also occur in
the arithmetic of classical signed root systems
(\citeproc{ref-ArdilaCastilloHenley2015}{Ardila et al. 2015, sec. 4.2}),
although the ambient lattice and the arbitrary cross-weighting here
differ. Hanusa and Zaslavsky's complete-incidence result concerns the
least common multiple of all minors of a Kronecker-product matrix, not
the gcd of maximal minors or Smith module of the present list
(\citeproc{ref-HanusaZaslavsky2011}{Hanusa and Zaslavsky 2011}). Graph
critical groups provide useful Smith-form context
(\citeproc{ref-Lorenzini2008}{Lorenzini 2008}), but the present cokernel
is not a reduced Laplacian cokernel and no chip-firing or
monodromy-pairing interpretation is asserted.

Arithmetic multiplicity is likewise established language. For a
represented list, \(m(E)\) is its saturation index and obeys the GCD
rule; over the complex torus it counts connected components of the
common character kernel (\citeproc{ref-Moci2012}{Moci 2012};
\citeproc{ref-DAdderioMoci2013}{D'Adderio and Moci 2013}). Matroids over
\(\mathbb Z\) and over DVRs retain the quotient module and local
invariant sequence rather than only its order
(\citeproc{ref-FinkMoci2016}{Fink and Moci 2016}). Accordingly, neither
\(h=m(E)\), the determinant gcd, nor the toric component interpretation
is claimed as new.

The expanded audit therefore changes the emphasis. Theorem 4.1 is an
explicit universal non-monic affine lift of classical square and monic
identities, and Corollary 5.1 is its formal Cauchy--Binet contraction.
Theorem 6.1 is a factorisation-specific cross-weighted application of
classical signed-incidence theory. The strongest candidate result is the
package formed by Theorems 7.1, 7.3, and 7.4 and Corollaries 7.2 and
7.5: the one-generator local module, every exact valuation, the
symmetric global formula, spanning-tree-gcd equality, and the scaled
Smith form for the complete cross-weighted list in the diagonal
quotient. The regular-unit interpretation is specific to the
factorisation open. Theorem 8.1 is best regarded as a scheme-theoretic
synthesis of classical labelled root partitions with the standard degree
of a torus isogeny.

Within the expanded bounded corpus, we did not locate the exact
all-factor affine Plücker tensor with its orientation, the local
presentation or closed index formula for this cross-weighted
diagonal-quotient list, or the tree-gcd theorem in this setting. This is
a bounded negative search result, not proof of priority. Equivalent
results may exist in determinant-complex, signed-graphic,
arithmetic-matroid, toric-arrangement, weighted-incidence, or
representation language. Independent specialist search and proof
reproduction remain necessary; no absolute priority claim is made.

\subsection{10.2 Mathematical
limitations}\label{mathematical-limitations}

Equation (7.1) classifies regular units only on the universal
pairwise-coprime open over \(\mathbb Z\), and on its polynomial-ring
analogue over a field. Under arbitrary base change, resultants may split
and the base may acquire new units. Arbitrary polynomial semi-invariants
are not classified, and they may become units after deleting additional
coefficient divisors.

The full unit-character index is always realised by a square
Laurent-unit basis. Saturation does not by itself give a basis made from
individual resultants or nonnegative resultant monomials. The example
\((1,2,2)\) has a nonnegative unimodular monomial basis, but the
corresponding semigroup-basis problem is open in general. Likewise, the
raw rectangular all-edge map has degree \([M:L]\) onto each
branch-labelled image; possible identifications of different branch
translates in an unlabelled rectangular target are not classified.

The numerical corank and pairwise-coprime regularity criterion for monic
multifactor multiplication are already described by Chaperon and López
de Medrano (\citeproc{ref-ChaperonLopezDeMedrano2009}{Chaperon and López
de Medrano 2009}, Theorem 3). What remains untreated here is therefore
not the set-theoretic rank classification, but the scheme-theoretic
structure of higher collision strata: their Fitting ideals,
multiplicities, embedded components, subresultant descriptions, and
behaviour under arbitrary base change. The maximal-minor ideal above
handles only the generic rank-loss divisor and its generic height-one
multiplicity.

Most importantly, local determinant control does not solve the global
affine-slice problem. A full-rank character matrix fixes the scaling
torus up to a finite group scheme and can make a normalised
factorisation chart étale off \(\Delta=0\). To obtain a polynomial
Keller map on affine space one must also prove that a selected level set
is isomorphic to affine space and that the restricted multiplication map
has the required global fibre or properness behaviour. No such
recognition theorem follows from the determinant line.

\subsection{10.3 Assurance limitations}\label{assurance-limitations}

The parent two-factor theorem is an immutable unrefereed research
candidate (\citeproc{ref-AnonymousBorderedJacobian2026}{Anonymous
2026}). The parent venue and this extension share a producer/research
process. The present induction has not been independently reproduced or
formally verified. The source search is bounded, the programs are
producer-written, and the manuscript has not undergone external peer
review. These limits apply even if every local build and exact check
passes.

\subsection{10.4 High-value next
problems}\label{high-value-next-problems}

The immediate next mathematical tasks are narrower than the programme's
largest aspirations.

\begin{enumerate}
\def\labelenumi{\arabic{enumi}.}
\tightlist
\item
  Obtain an independently written proof of the affine Plücker lift,
  including every orientation convention, and independently reproduce
  the local Smith presentation and global formula.
\item
  Formalise the determinant-line composition lemma, the local Tietze
  elimination, and the Fitting-ideal consequences in a proof assistant.
\item
  Classify when the unit-character lattice has a basis of nonnegative
  resultant monomials, then investigate arbitrary semi-invariant
  vertical Jacobians after additional divisor deletions.
\item
  Refine the classical numerical corank formula to scheme-theoretic
  higher collision strata using subresultants and Fitting ideals.
\item
  Determine whether the complete unit-character group controls monodromy
  or compactification data beyond the generic torus chart.
\item
  Treat affine-slice recognition as a separate global geometry problem
  before drawing Keller or Hessian conclusions.
\end{enumerate}

The most credible programme-level target is a classification of
factorisation-derived étale charts and the isolation, if true, of
exceptional affine-space slices. A four-dimensional Hessian application
remains a conditional long-range problem, not a consequence suggested by
the present proof alone.

\section{11. Conclusion}\label{conclusion}

The strongest candidate calculation in this paper is now exact. Theorem
7.1 reduces the complete primitive resultant-character quotient at every
prime to one generator with the relations \(D_a\) and \(2e_ie_j\).
Corollary 7.2 gives every valuation of its order; Theorem 7.3 assembles
those valuations into a symmetric global formula; and Theorem 7.4 shows
that spanning-tree minors already determine the full arithmetic-matroid
multiplicity. Scaling then gives the complete nonprimitive Smith form
and the connected/étale decomposition of the unit-character residual
group.

The route to this calculation begins with a classical multiplication
Jacobian. Theorem 4.1 packages that antecedent as the exact affine
determinant-line identity

\[
 \text{maximal-minor tensor}
 =\text{pairwise resultant divisor}
 \times\text{torus Plücker tensor},
\]

with an explicit integral orientation. Multi-bordering pairs the torus
tensor with the vertical Jacobian of the chosen normalising functions.
For semi-invariants, this pairing is an integer character determinant.
Pairwise resultants turn the character calculation into cross-weighted
signed-incidence theory, from which the graph-minor formula follows.
Standard arithmetic-matroid language then identifies the complete-list
gcd as \(m(E)\), while the local presentation evaluates that standard
invariant and its full quotient module for this factorisation-derived
list.

The rank-one degree difference is therefore the rank-one character
determinant inside a broader Smith theory. In higher torus rank, a
selected square character determinant controls a particular finite
chart, whereas the complete unit-character lattice controls the residual
group intrinsic to the fixed pairwise-coprime open. Theorem 8.1
translates square Smith data into total, separable, and inseparable
scheme degrees through a finite-locally-free factorisation. These
distinctions are the durable conclusion. The publication boundary
remains provisional until the proofs and strengthened literature
position receive genuinely independent specialist review.

\section{Data and code availability}\label{data-and-code-availability}

The research package contains the manuscript source, bibliography, exact
SymPy and FLINT verification programs, the integer character-lattice
program, normal and optimized receipts, the Stage-4 local-Smith verifier
and receipts, a manifest, the proof and literature audit notes, and a
formalisation blueprint. The package is a local research-stage child
artefact at the time of writing. No public repository, archival deposit,
DOI, or independent reproduction for this extension is asserted.

\section{Submission metadata and
licence}\label{submission-metadata-and-licence}

The author name, affiliation, complete contribution statement, final
conflict declaration, funding statement, and distribution licence have
not been supplied. They are not inferred in this research draft. The
artefact is not submission-ready or release-ready until the responsible
author supplies and approves those fields.

\section{Ethics declaration}\label{ethics-declaration}

This work is theoretical mathematics and involved no human participants,
personal data, animals, clinical intervention, or sensitive field data.
The principal research-ethics risks are attribution error, novelty
overstatement, and inflation of producer checks into independent
verification; the manuscript states explicit controls for each.

\section{Author contributions}\label{author-contributions}

\textbf{Conceptualization:} {[}to be supplied{]}. \textbf{Formal
analysis:} {[}to be supplied{]}. \textbf{Methodology:} {[}to be
supplied{]}. \textbf{Software:} {[}to be supplied{]}.
\textbf{Validation:} {[}to be supplied{]}. \textbf{Investigation:} {[}to
be supplied{]}. \textbf{Writing -- original draft:} {[}to be
supplied{]}. \textbf{Writing -- review and editing:} {[}to be
supplied{]}.

\section{Conflict of interest}\label{conflict-of-interest}

The author information has not yet been supplied. The parent Evidence
Press candidate and this extension share a producer/research process;
Evidence Press availability is not independent scholarly endorsement. No
other conflict is asserted at this stage.

\section{Funding}\label{funding}

Funding information has not been supplied. The manuscript makes no claim
of external financial support.

\section{AI-use disclosure}\label{ai-use-disclosure}

AI systems, including OpenAI language-model tooling, assisted source
discovery, proof exploration, exact-check implementation, drafting, and
internal review. The mathematical statements, citations, code outputs,
and assurance labels remain the responsibility of the named author or
authors once supplied. AI assistance and producer-side computation do
not constitute independent verification. A public version should retain
a venue-appropriate disclosure with the specific systems and human
review responsibilities identified.

\section{Appendix A. Orientation and indexing
conventions}\label{appendix-a.-orientation-and-indexing-conventions}

\begin{enumerate}
\def\labelenumi{\arabic{enumi}.}
\tightlist
\item
  Coefficients are ordered by factor, then from \(X^{d_i}\) to
  \(Y^{d_i}\).
\item
  Sylvester rows are shifts of the earlier factor followed by shifts of
  the later factor.
\item
  Deleted source columns use zero-based indices.
\item
  The Plücker sign is \((-1)^{\sum I}\).
\item
  The torus basis is \(e_b-e_k\), ordered by \(b=1,\ldots,k-1\).
\item
  Gradient rows follow multiplication rows and are ordered
  \(g_1,\ldots,g_{k-1}\).
\item
  The multi-border Laplace sign is \((-1)^{nt+\binom t2}\).
\end{enumerate}

Changing any convention changes intermediate signs. The invariant
content is the determinant-line equality after orientations are
transformed together.

\section{Appendix B. Verification
boundary}\label{appendix-b.-verification-boundary}

The default receipts verify the exact cases listed in Section 9 and
include deliberate negative controls. Larger exploratory cases were also
tested but are not part of the frozen default receipt. No finite list of
symbolic cases proves the universal theorem. The proof is the induction
in Sections 3 and 4; the programs are regression tests for its
coordinate consequences.

\section{Appendix C. Formalisation
outline}\label{appendix-c.-formalisation-outline}

A proof-assistant kernel can be divided into the following layers:

\begin{enumerate}
\def\labelenumi{\arabic{enumi}.}
\tightlist
\item
  coefficient convolution and the universal multiplication matrix;
\item
  the self-contained two-factor complementary-minor identity of Lemma
  E.1;
\item
  signed complementary minors and the block-composition lemma;
\item
  resultant multiplicativity and the all-factor induction;
\item
  Laplace expansion, Cauchy--Binet, and multi-border contraction;
\item
  integer character matrices and Smith determinantal divisors; and
\item
  the weighted incidence reduction.
\end{enumerate}

The universal polynomial identities should be formalised before
geometric language about divisors or étaleness. This order keeps the
trusted kernel small and isolates the sign-sensitive finite-dimensional
linear algebra.

\section{Appendix D. Bounded antecedent and claim
matrix}\label{appendix-d.-bounded-antecedent-and-claim-matrix}

The following matrix records the source-to-claim boundary used in this
draft. It is a bounded map of the examined corpus, not a proof that no
closer antecedent exists.

{\def\LTcaptype{none} % do not increment counter
\begin{longtable}[]{@{}
  >{\raggedright\arraybackslash}p{(\linewidth - 6\tabcolsep) * \real{0.2500}}
  >{\raggedright\arraybackslash}p{(\linewidth - 6\tabcolsep) * \real{0.2500}}
  >{\raggedright\arraybackslash}p{(\linewidth - 6\tabcolsep) * \real{0.2500}}
  >{\raggedright\arraybackslash}p{(\linewidth - 6\tabcolsep) * \real{0.2500}}@{}}
\toprule\noalign{}
\begin{minipage}[b]{\linewidth}\raggedright
Object in this paper
\end{minipage} & \begin{minipage}[b]{\linewidth}\raggedright
Closest examined antecedent
\end{minipage} & \begin{minipage}[b]{\linewidth}\raggedright
What the antecedent establishes
\end{minipage} & \begin{minipage}[b]{\linewidth}\raggedright
Candidate increment evaluated here
\end{minipage} \\
\midrule\noalign{}
\endhead
\bottomrule\noalign{}
\endlastfoot
Square two-factor product Jacobian & Artin, Basu--Pollack--Roy, and
Bhargava--Cremona--Fisher--Gajović (\citeproc{ref-Artin2022}{Artin
2022}; \citeproc{ref-BasuPollackRoy2006}{Basu et al. 2006};
\citeproc{ref-BhargavaCremonaFisherGajovic2022}{Bhargava et al. 2022}) &
The coefficient-product Jacobian is a Sylvester matrix, hence the
resultant, including a monic-times-arbitrary identity over an arbitrary
ring & Retain every affine coefficient direction and record the complete
complementary-minor vector with a fixed integral orientation \\
Projective and monic multifactor rank & Chipalkatti and Chaperon--López
de Medrano (\citeproc{ref-Chipalkatti2003}{Chipalkatti 2003};
\citeproc{ref-ChaperonLopezDeMedrano2009}{Chaperon and López de Medrano
2009}) & Projective tangent injectivity for coprime factors, and the
monic arbitrary-factor corank and pairwise-coprimality criterion &
Refine the generic rank-loss set to an exact universal maximal-minor
ideal with generic height-one multiplicities \\
Scalar product \(\prod_{i<j}R_{ij}\) & Mahatab--Sampath
(\citeproc{ref-MahatabSampath2015}{Mahatab and Sampath 2015}) &
Determinant of the polynomial Chinese-remainder map for monic factors &
Restore all affine leading-coefficient directions and identify the full
complementary-minor tensor \\
Resultant divisibility in determinantal systems & Chardin
(\citeproc{ref-Chardin1993}{Chardin 1993}) & Resultant factors and
generic gcd behaviour for generalized Sylvester/Koszul minors & Give
every signed maximal minor explicitly for the coefficient differential
of labelled multiplication \\
Determinant-of-complexes framework & Jouanolou, Demazure, and GKZ
(\citeproc{ref-Jouanolou1991}{Jouanolou 1991},
\citeproc{ref-Jouanolou1997}{1997}; \citeproc{ref-Demazure2012}{Demazure
2012}; \citeproc{ref-GelfandKapranovZelevinsky1994}{Gelfand et al.
1994}) & General resultant constructions, inertia forms, and
discriminantal determinants & Isolate the torus-kernel Plücker line and
its border contraction in this affine factorisation map \\
Two-factor affine identity & Parent candidate
(\citeproc{ref-AnonymousBorderedJacobian2026}{Anonymous 2026}) &
Rank-one complementary-minor and bordered identities & Prove the
all-factor induction, multi-border formula, and higher-rank character
interpretation \\
Projective multiplication fibres & Breiding--Kohn--Sturmfels
(\citeproc{ref-BreidingKohnSturmfels2024}{Breiding et al. 2024}) &
Projective polynomial multiplication and finite root-partition geometry
& Give an fpqc-local finite-locally-free factorisation that composes the
root-partition cover with a translated torus isogeny \\
Linear-factor quotient geometry & Kurth (\citeproc{ref-Kurth1997}{Kurth
1997}) & Ordered linear factors modulo scaling and permutation actions &
Retain arbitrary positive degrees and compute the resultant-character
lattice integrally \\
Signed support of character bases & Zaslavsky with erratum,
Grossman--Kulkarni--Schochetman, and Hessert--Mallik
(\citeproc{ref-Zaslavsky1982}{Zaslavsky 1982},
\citeproc{ref-Zaslavsky1983Erratum}{1983};
\citeproc{ref-GrossmanKulkarniSchochetman1995}{Grossman et al. 1995};
\citeproc{ref-HessertMallik2023}{Hessert and Mallik 2023}) &
All-negative signed-incidence rank, dependence, odd-cycle support,
powers of two, and odd-unicyclic invertibility & Evaluate the appended
degree balance and cross-degree product and identify the characters with
pairwise resultants \\
Multiplicity of the complete character list & Moci and D'Adderio--Moci
(\citeproc{ref-Moci2012}{Moci 2012};
\citeproc{ref-DAdderioMoci2013}{D'Adderio and Moci 2013}) & Saturation
multiplicity, GCD rule, and complex toric component count & Give a
closed evaluation for the complete cross-weighted diagonal-quotient list
and its factorisation-open interpretation \\
Exact full-set quotient module & Fink--Moci; Ardila--Castillo--Henley as
an adjacent equal-weight root-list case
(\citeproc{ref-FinkMoci2016}{Fink and Moci 2016};
\citeproc{ref-ArdilaCastilloHenley2015}{Ardila et al. 2015}) &
Matroid-over-ring/DVR framework and arithmetic multiplicities for
classical signed-root lists & Prove the one-generator local
presentation, every valuation, the scaled global Smith form, and
connected/étale residual group factors \\
Complete-graph incidence subdeterminants & Hanusa--Zaslavsky
(\citeproc{ref-HanusaZaslavsky2011}{Hanusa and Zaslavsky 2011}) & Least
common multiple of all minors of a fixed Kronecker-product
complete-incidence matrix & Compute a gcd of maximal minors and the full
Smith module for a different cross-weighted signless list in the
diagonal quotient \\
Jacobian ideals of binary discriminants & D'Andrea--Chipalkatti
(\citeproc{ref-DAndreaChipalkatti2007}{D'Andrea and Chipalkatti 2007}) &
Structure and resolutions for coefficient-gradient ideals of binary
discriminants and resultants & Study the rectangular differential of
factor multiplication rather than the gradient ideal of the resultant
hypersurface \\
\end{longtable}
}

The affine Plücker identity is best viewed as a derived,
coordinate-complete refinement of the exact square antecedents. The
strongest conditional originality claim supported by this matrix is
narrower: the exact local module, closed valuation/global formula, and
tree-gcd equality for the cross-weighted complete-edge list were not
located in the bounded search. Specialist review could still identify an
equivalent formulation in weighted-incidence, arithmetic-matroid, toric,
logarithmic, or representation-theoretic language.

\section{Appendix E. Self-contained affine lift of the classical
two-factor Jacobian
identity}\label{appendix-e.-self-contained-affine-lift-of-the-classical-two-factor-jacobian-identity}

Classical sources identify a square two-factor product Jacobian with the
resultant (\citeproc{ref-Artin2022}{Artin 2022, sec. 1.8}, Lemma 1.8.5;
\citeproc{ref-BasuPollackRoy2006}{Basu et al. 2006, sec. 4.2.1};
\citeproc{ref-BhargavaCremonaFisherGajovic2022}{Bhargava et al. 2022},
Lemma 2.6). This appendix supplies the coordinate-complete affine
complementary-minor lift and fixes its sign in the conventions of this
paper. It also reproduces the parent rank-one base theorem
(\citeproc{ref-AnonymousBorderedJacobian2026}{Anonymous 2026}), so
Theorem 4.1 remains self-contained without asserting priority for the
underlying square Jacobian identity.

\textbf{Lemma E.1 (two-factor complementary minors over \(\mathbb Z\)).}
Let

\[
 A=\sum_{i=0}^{r}a_iX^{r-i}Y^i,\qquad
 B=\sum_{j=0}^{s}b_jX^{s-j}Y^j,
 \qquad r,s\ge1,
\]

and write \(AB=\sum_{h=0}^{r+s}c_hX^{r+s-h}Y^h\). Order source
coordinates as \(a_0,\ldots,a_r,b_0,\ldots,b_s\) and target coordinates
as \(c_0,\ldots,c_{r+s}\). Let \(M\) be the coefficient differential,

\[
 M[h,i]=b_{h-i},\qquad
 M[h,r+1+j]=a_{h-j},
\]

with out-of-range subscripts equal to zero, and put

\[
 \kappa=(a_0,\ldots,a_r,-b_0,\ldots,-b_s).
\]

Then

\[
 \boxed{
 (-1)^\ell\det M_{\widehat\ell}
 =(-1)^{r(s+1)}\operatorname{Res}(A,B)\kappa_\ell,
 \qquad 0\le\ell\le r+s+1.
 }
\tag{E.1}
\]

\textbf{Proof.} Both sides are polynomials with integer coefficients, so
it suffices to prove their equality on a nonempty Zariski-open subset
over \(\mathbb C\). Dehomogenise by \(t=Y/X\), work where
\(a_rb_s\ne0\), and write

\[
 a(t)=a_r\prod_{i=1}^{r}(t-\alpha_i),\qquad
 b(t)=b_s\prod_{j=1}^{s}(t-\beta_j),
\]

with all \(\alpha_i,\beta_j\) pairwise distinct. In the Sylvester
convention of Section 2,

\[
 \operatorname{Vand}(x_1,\ldots,x_m)
 :=\prod_{1\le u<v\le m}(x_v-x_u).
\]

Every Vandermonde below uses the displayed order of its arguments. With
this orientation,

\[
 \operatorname{Res}(A,B)
 =a_r^sb_s^r\prod_{i,j}(\beta_j-\alpha_i)
 =b_s^r\prod_j a(\beta_j)
 =(-1)^{rs}a_r^s\prod_i b(\alpha_i).
\tag{E.2}
\]

For completeness, transpose the Sylvester matrix and view it as
\((u,w)\mapsto ua+wb\). Evaluating the output at
\((\alpha_1,\ldots,\alpha_r,\beta_1,\ldots,\beta_s)\) makes the matrix
block anti-diagonal. Its determinant is

\[
 (-1)^{rs}
 \left(\prod_i b(\alpha_i)\right)\operatorname{Vand}(\alpha)
 \left(\prod_j a(\beta_j)\right)\operatorname{Vand}(\beta).
\]

Division by the evaluation determinant

\[
 \operatorname{Vand}(\alpha)\operatorname{Vand}(\beta)
 \prod_{i,j}(\beta_j-\alpha_i)
\]

gives (E.2); the specialisation \((A,B)=(X^r,Y^s)\) fixes the
normalisation as \(+1\).

Choose \(\tau\in\mathbb C\) distinct from every root and evaluate the
degree-\((r+s)\) target at

\[
 \alpha_1,\ldots,\alpha_r,\beta_1,\ldots,\beta_s,\tau.
\]

Let \(V\) be this evaluation matrix. Its determinant is the Vandermonde
in the displayed order. At \(\alpha_i\), the evaluated Jacobian row is

\[
 b(\alpha_i)(1,\alpha_i,\ldots,\alpha_i^r\mid0);
\]

at \(\beta_j\), it is

\[
 a(\beta_j)(0\mid1,\beta_j,\ldots,\beta_j^s).
\]

The \(\tau\)-row is the sum of its \(A\)-block and \(B\)-block parts.
After one source column is deleted, one summand has deficient block rank
and vanishes.

We also need the gap-Vandermonde identity. For \(m\) variables
\(x_1,\ldots,x_m\), let \(W_{\widehat h}(x)\) be the \(m\)-square
evaluation matrix with exponents \(\{0,\ldots,m\}\setminus\{h\}\).
Comparing coefficients of \(T^h\) in the Vandermonde determinant on
\((x_1,\ldots,x_m,T)\) gives

\[
 \det W_{\widehat h}(x)
 =e_{m-h}(x)\operatorname{Vand}(x).
\tag{E.3}
\]

Suppose first that the deleted column is \(b_j\), of global index
\(\ell=r+1+j\). The surviving term places the \(\tau\)-row in the
\(A\)-block. Moving it past the \(s\) beta rows contributes \((-1)^s\),
and block expansion gives

\[
 \begin{aligned}
 \det(VM_{\widehat\ell})
 ={}&(-1)^s
 \left(\prod_i b(\alpha_i)\right)b(\tau)
 \operatorname{Vand}(\alpha,\tau)\\
 &\quad\cdot
 \left(\prod_j a(\beta_j)\right)
 \det W_{\widehat j}(\beta).
 \end{aligned}
\tag{E.4}
\]

Use (E.2)--(E.3),

\[
 b(\tau)=b_s\prod_j(\tau-\beta_j),\qquad
 \prod_i b(\alpha_i)=(-1)^{rs}b_s^r
       \prod_{i,j}(\beta_j-\alpha_i),
\]

and divide by \(\det V\). All Vandermonde and \(\tau\)-factors cancel,
leaving

\[
 \det M_{\widehat\ell}
 =(-1)^{s(r+1)}b_se_{s-j}(\beta)\operatorname{Res}(A,B).
\tag{E.5}
\]

Since \(b_j=(-1)^{s-j}b_se_{s-j}(\beta)\),

\[
 \begin{aligned}
 (-1)^\ell\det M_{\widehat\ell}
 &=(-1)^{r+1+j+s(r+1)+s-j}
   b_j\operatorname{Res}(A,B)\\
 &=(-1)^{r(s+1)}(-b_j)\operatorname{Res}(A,B),
 \end{aligned}
\]

which is (E.1) in the \(B\)-block.

If the deleted column is \(a_i\), its global index is \(\ell=i\). The
surviving term places the \(\tau\)-row in the \(B\)-block, and the
existing row order is already block compatible:

\[
 \begin{aligned}
 \det(VM_{\widehat i})
 ={}&\left(\prod_i b(\alpha_i)\right)
 \det W_{\widehat i}(\alpha)\\
 &\quad\cdot
 \left(\prod_j a(\beta_j)\right)a(\tau)
 \operatorname{Vand}(\beta,\tau).
 \end{aligned}
\tag{E.6}
\]

The same cancellations yield

\[
 \det M_{\widehat i}
 =(-1)^{rs}a_re_{r-i}(\alpha)\operatorname{Res}(A,B)
 =(-1)^{rs+r-i}a_i\operatorname{Res}(A,B).
\tag{E.7}
\]

Multiplication by \((-1)^i\) gives (E.1) in the \(A\)-block. The
identity holds on the chosen dense open subset, hence as a polynomial
identity over \(\mathbb Z\). \(\square\)

For reference, the two deleted-column parity calculations are:

{\def\LTcaptype{none} % do not increment counter
\begin{longtable}[]{@{}
  >{\raggedright\arraybackslash}p{(\linewidth - 10\tabcolsep) * \real{0.1304}}
  >{\raggedleft\arraybackslash}p{(\linewidth - 10\tabcolsep) * \real{0.1739}}
  >{\raggedleft\arraybackslash}p{(\linewidth - 10\tabcolsep) * \real{0.1739}}
  >{\raggedleft\arraybackslash}p{(\linewidth - 10\tabcolsep) * \real{0.1739}}
  >{\raggedleft\arraybackslash}p{(\linewidth - 10\tabcolsep) * \real{0.1739}}
  >{\raggedleft\arraybackslash}p{(\linewidth - 10\tabcolsep) * \real{0.1739}}@{}}
\toprule\noalign{}
\begin{minipage}[b]{\linewidth}\raggedright
Deleted block
\end{minipage} & \begin{minipage}[b]{\linewidth}\raggedleft
Global index
\end{minipage} & \begin{minipage}[b]{\linewidth}\raggedleft
Sign of \(\det M_{\widehat\ell}\) after coefficient substitution
\end{minipage} & \begin{minipage}[b]{\linewidth}\raggedleft
Cofactor sign
\end{minipage} & \begin{minipage}[b]{\linewidth}\raggedleft
Kernel-coordinate sign
\end{minipage} & \begin{minipage}[b]{\linewidth}\raggedleft
Remaining scalar sign
\end{minipage} \\
\midrule\noalign{}
\endhead
\bottomrule\noalign{}
\endlastfoot
\(A\), column \(a_i\) & \(i\) & \(rs+r-i\) & \(i\) & \(0\) for \(a_i\) &
\(r(s+1)\) \\
\(B\), column \(b_j\) & \(r+1+j\) & \(s(r+1)+s-j\) & \(r+1+j\) & \(1\)
for \(-b_j\) & \(r(s+1)\) \\
\end{longtable}
}

All entries in the last four columns are exponents of \((-1)\), read
modulo two. This table records the cancellations between (E.4)--(E.5)
and (E.6)--(E.7) without changing the proof.

Expansion along an arbitrary final row \(v\) now gives the equivalent
bordered form

\[
 \det\begin{pmatrix}M\\v\end{pmatrix}
 =(-1)^{s(r+1)+1}\operatorname{Res}(A,B)
   \langle v,\kappa\rangle,
\]

which also confirms the \(k=2\) specialisation of Corollary 5.1.

\section*{References}\label{references}
\addcontentsline{toc}{section}{References}

\protect\phantomsection\label{refs}
\begin{CSLReferences}{1}{1}
\bibitem[\citeproctext]{ref-AnonymousBorderedJacobian2026}
Anonymous. 2026. \emph{The Bordered Jacobian of Binary-Form
Multiplication: An Integral Foundation for the Degree-Difference
Principle}. Evidence Press research candidate deposited in Zenodo.
\url{https://doi.org/10.5281/zenodo.21855302}.

\bibitem[\citeproctext]{ref-ArdilaCastilloHenley2015}
Ardila, Federico, Federico Castillo, and Michael Henley. 2015. {``The
Arithmetic {T}utte Polynomials of the Classical Root Systems.''}
\emph{International Mathematics Research Notices} 2015 (12): 3830--77.
\url{https://doi.org/10.1093/imrn/rnu050}.

\bibitem[\citeproctext]{ref-Artin2022}
Artin, Michael. 2022. \emph{Algebraic Geometry: Notes on a Course}. Vol.
222. Graduate Studies in Mathematics. American Mathematical Society.
\url{https://doi.org/10.1090/gsm/222}.

\bibitem[\citeproctext]{ref-BasuPollackRoy2006}
Basu, Saugata, Richard Pollack, and Marie-Françoise Roy. 2006.
\emph{Algorithms in Real Algebraic Geometry}. 2nd ed. Vol. 10.
Algorithms and Computation in Mathematics. Springer.
\url{https://doi.org/10.1007/3-540-33099-2}.

\bibitem[\citeproctext]{ref-BhargavaCremonaFisherGajovic2022}
Bhargava, Manjul, John Cremona, Tom Fisher, and Stevan Gajović. 2022.
{``The Density of Polynomials of Degree {\(n\)} over {\(\mathbb Z_p\)}
Having Exactly {\(r\)} Roots in {\(\mathbb Q_p\)}.''} \emph{Proceedings
of the London Mathematical Society} 124 (5): 713--36.
\url{https://doi.org/10.1112/plms.12438}.

\bibitem[\citeproctext]{ref-BreidingKohnSturmfels2024}
Breiding, Paul, Kathlén Kohn, and Bernd Sturmfels. 2024. \emph{Metric
Algebraic Geometry}. Vol. 53. Oberwolfach Seminars. Birkh{ä}user.
\url{https://doi.org/10.1007/978-3-031-51462-3}.

\bibitem[\citeproctext]{ref-Cayley1848}
Cayley, Arthur. 1848. {``On the Theory of Elimination.''} \emph{The
Cambridge and Dublin Mathematical Journal} 3: 116--20.
\url{https://www.deutsche-digitale-bibliothek.de/item/BWTZ6SV3BZSX4RJUAWQFVZ6R277IQRIY}.

\bibitem[\citeproctext]{ref-ChaperonLopezDeMedrano2009}
Chaperon, Marc, and Santiago López de Medrano. 2009. {``Some
Regularities and Singularities Appearing in the Study of Polynomials and
Operators.''} \emph{Ast{é}risque}, no. 323: 123--60.
\url{https://www.numdam.org/item/AST_2009__323__123_0/}.

\bibitem[\citeproctext]{ref-Chardin1993}
Chardin, Marc. 1993. {``The Resultant via a {Koszul} Complex.''} In
\emph{Computational Algebraic Geometry}, edited by Frédéric Eyssette and
André Galligo, vol. 109. Progress in Mathematics. Birkh{ä}user Boston.
\url{https://doi.org/10.1007/978-1-4612-2752-6_3}.

\bibitem[\citeproctext]{ref-Chipalkatti2003}
Chipalkatti, Jaydeep V. 2003. {``On Equations Defining Coincident Root
Loci.''} \emph{Journal of Algebra} 267 (1): 246--71.
\url{https://doi.org/10.1016/S0021-8693(03)00336-3}.

\bibitem[\citeproctext]{ref-DAdderioMoci2013}
D'Adderio, Michele, and Luca Moci. 2013. {``Arithmetic Matroids, the
{T}utte Polynomial and Toric Arrangements.''} \emph{Advances in
Mathematics} 232 (1): 335--67.
\url{https://doi.org/10.1016/j.aim.2012.09.001}.

\bibitem[\citeproctext]{ref-DAndreaChipalkatti2007}
D'Andrea, Carlos, and Jaydeep Chipalkatti. 2007. {``On the Jacobian
Ideal of the Binary Discriminant.''} \emph{Collectanea Mathematica} 58
(2): 155--80. \url{https://arxiv.org/abs/math/0601705}.

\bibitem[\citeproctext]{ref-Demazure2012}
Demazure, Michel. 2012. {``R{é}sultant, Discriminant.''}
\emph{L'Enseignement Math{é}matique} 58 (3/4): 333--73.
\url{https://doi.org/10.4171/LEM/58-3-5}.

\bibitem[\citeproctext]{ref-FinkMoci2016}
Fink, Alex, and Luca Moci. 2016. {``Matroids over a Ring.''}
\emph{Journal of the European Mathematical Society} 18 (4): 681--731.
\url{https://doi.org/10.4171/JEMS/600}.

\bibitem[\citeproctext]{ref-GelfandKapranovZelevinsky1994}
Gelfand, Israel M., Mikhail M. Kapranov, and Andrei V. Zelevinsky. 1994.
\emph{Discriminants, Resultants, and Multidimensional Determinants}.
Mathematics: Theory and Applications. Birkh{ä}user.
\url{https://doi.org/10.1007/978-0-8176-4771-1}.

\bibitem[\citeproctext]{ref-GrossmanKulkarniSchochetman1995}
Grossman, Jerrold W., Devadatta M. Kulkarni, and Irwin E. Schochetman.
1995. {``On the Minors of an Incidence Matrix and Its {Smith} Normal
Form.''} \emph{Linear Algebra and Its Applications} 218: 213--24.
\url{https://doi.org/10.1016/0024-3795(93)00173-W}.

\bibitem[\citeproctext]{ref-HanusaZaslavsky2011}
Hanusa, Christopher R. H., and Thomas Zaslavsky. 2011. {``Determinants
in the {K}ronecker Product of Matrices: The Incidence Matrix of a
Complete Graph.''} \emph{Linear and Multilinear Algebra} 59 (4):
399--411. \url{https://doi.org/10.1080/03081081003586852}.

\bibitem[\citeproctext]{ref-HessertMallik2023}
Hessert, Ryan, and Sudipta Mallik. 2023. {``The Inverse of the Incidence
Matrix of a Unicyclic Graph.''} \emph{Linear and Multilinear Algebra} 71
(4): 513--27. \url{https://doi.org/10.1080/03081087.2022.2035307}.

\bibitem[\citeproctext]{ref-Jouanolou1991}
Jouanolou, Jean-Pierre. 1991. {``Le Formalisme Du r{é}sultant.''}
\emph{Advances in Mathematics} 90 (2): 117--263.
\url{https://doi.org/10.1016/0001-8708(91)90031-2}.

\bibitem[\citeproctext]{ref-Jouanolou1995}
Jouanolou, Jean-Pierre. 1995. {``Aspects Invariants de
l'{é}limination.''} \emph{Advances in Mathematics} 114 (1): 1--174.
\url{https://doi.org/10.1006/aima.1995.1042}.

\bibitem[\citeproctext]{ref-Jouanolou1997}
Jouanolou, Jean-Pierre. 1997. {``Formes d'inertie Et r{é}sultant: Un
Formulaire.''} \emph{Advances in Mathematics} 126 (2): 119--250.
\url{https://doi.org/10.1006/aima.1996.1609}.

\bibitem[\citeproctext]{ref-Kurth1997}
Kurth, Alexandre. 1997. {``{\(\mathrm{SL}_2\)}-Equivariant Polynomial
Automorphisms of the Binary Forms.''} \emph{Annales de l'Institut
Fourier} 47 (2): 585--97. \url{https://doi.org/10.5802/aif.1574}.

\bibitem[\citeproctext]{ref-Lorenzini2008}
Lorenzini, Dino J. 2008. {``Smith Normal Form and Laplacians.''}
\emph{Journal of Combinatorial Theory, Series B} 98 (6): 1271--300.
\url{https://doi.org/10.1016/j.jctb.2008.02.002}.

\bibitem[\citeproctext]{ref-MahatabSampath2015}
Mahatab, Kamalakshya, and Kannappan Sampath. 2015. {``Chinese Remainder
Theorem for Cyclotomic Polynomials in {\(\mathbf{Z}[X]\)}.''}
\emph{Journal of Algebra} 435: 223--62.
\url{https://doi.org/10.1016/j.jalgebra.2015.04.006}.

\bibitem[\citeproctext]{ref-Moci2012}
Moci, Luca. 2012. {``A {T}utte Polynomial for Toric Arrangements.''}
\emph{Transactions of the American Mathematical Society} 364 (2):
1067--88. \url{https://doi.org/10.1090/S0002-9947-2011-05491-7}.

\bibitem[\citeproctext]{ref-StacksProject2026}
The Stacks Project Authors. 2026. \emph{The {Stacks} Project}. Online
reference. \url{https://stacks.math.columbia.edu}.

\bibitem[\citeproctext]{ref-Zaslavsky1982}
Zaslavsky, Thomas. 1982. {``Signed Graphs.''} \emph{Discrete Applied
Mathematics} 4 (1): 47--74.
\url{https://doi.org/10.1016/0166-218X(82)90033-6}.

\bibitem[\citeproctext]{ref-Zaslavsky1983Erratum}
Zaslavsky, Thomas. 1983. {``Signed Graphs {[}Erratum{]}.''}
\emph{Discrete Applied Mathematics} 5 (2): 248.
\url{https://doi.org/10.1016/0166-218X(83)90047-1}.

\end{CSLReferences}

\end{document}
